% FORGE, ICLR 2027 conference manuscript.
% Quantitative inputs are byte-identical snapshots of the recorded source artifacts.
% See iclr_provenance.json before refreshing generated tables, figures, or bibliography files.

\documentclass{article}
\usepackage{iclr2027_conference,times}
% Anonymous review mode; enable \iclrfinalcopy only for an accepted camera-ready version.

\usepackage[utf8]{inputenc}
\usepackage[T1]{fontenc}
\usepackage{nicefrac}
\usepackage{microtype}
\usepackage{amsmath,amssymb,amsthm,mathrsfs}
\usepackage{algorithm}
\usepackage{algpseudocode}
\usepackage{graphicx}
\usepackage{pgfplots}
\usepackage{pgfplotstable}
\pgfplotsset{compat=1.18}
\usepackage{tikz}
\usetikzlibrary{arrows.meta,calc,positioning,patterns,shapes.arrows,fit,backgrounds}
\usepackage{array}
\usepackage{booktabs}
\usepackage{colortbl}
\usepackage{longtable}
\usepackage[most]{tcolorbox}
\usepackage{float}
\usepackage{placeins}
\usepackage{chemfig}
\usepackage{xcolor}
\usepackage{url}
\usepackage[hidelinks]{hyperref}

% Figures and generated rows are local, hash-recorded copies of the source manuscript inputs.
\graphicspath{{./}}

\newtheorem{theorem}{Theorem}[section]
\newtheorem{proposition}[theorem]{Proposition}
\newtheorem{lemma}[theorem]{Lemma}
\newtheorem{corollary}[theorem]{Corollary}
\theoremstyle{definition}
\newtheorem{definition}[theorem]{Definition}
\newtheorem{assumption}[theorem]{Assumption}
\newtheorem{invariant}[theorem]{Implementation invariant}

% Formal statements use a quiet lavender field to separate assumptions and guarantees from the
% surrounding argument. The box may break across pages; proofs intentionally remain unboxed.
\definecolor{forgeformalbg}{HTML}{F0EFFF}
\definecolor{forgerow}{HTML}{E9E7FB}
\tcbset{
  forge formal block/.style={
    enhanced jigsaw,
    breakable,
    colback=forgeformalbg,
    colframe=forgeformalbg,
    boxrule=0pt,
    arc=0pt,
    outer arc=0pt,
    boxsep=0pt,
    left=8pt,
    right=8pt,
    top=7pt,
    bottom=7pt,
    before skip=8pt,
    after skip=8pt
  }
}
\tcolorboxenvironment{definition}{forge formal block}
\tcolorboxenvironment{assumption}{forge formal block}
\tcolorboxenvironment{proposition}{forge formal block}
\tcolorboxenvironment{theorem}{forge formal block}
\tcolorboxenvironment{lemma}{forge formal block}
\tcolorboxenvironment{corollary}{forge formal block}
\tcolorboxenvironment{invariant}{forge formal block}

% Unresolved experimental results remain visually explicit until a verified artifact supplies
% the value. Never replace one of these placeholders from an unverified log or planning document.
\newcommand{\RESULTPENDING}[1]{\textbf{\textcolor{red}{[PENDING: #1]}}}
% Generated from the hash-pinned completed-evidence contract. Do not edit numerical values here.
\input{sections/result_placeholders.tex}
\input{generated/completed_evidence_macros.tex}
% Superseding Table 1 values from the matched final conditioned, null and cyclic evaluations.
\input{generated/iclr_table1_macros.tex}
% Superseding Table 3 values from the final graph Transformer population and frozen route index.
\input{generated/iclr_table3_route_macros.tex}
% Superseding Table 10 values from the final method-blind common-Ugi route-evidence adjudication.
\input{generated/iclr_table10_macros.tex}
% Every computed value cited in prose is regenerated from the same final evidence as its table.
\input{generated/iclr_prose_results_macros.tex}

% Appendix generated-structure atlas. The hash-pinned row file carries structures and SMILES only
% and defers every length, rule and type style to the definitions below, so the appendix can be
% retuned without reissuing a provenance artifact.
\definecolor{atlasrule}{gray}{0.74}
\definecolor{atlasink}{gray}{0.42}
\newlength{\atlasgraphwidth}
\newlength{\atlasprecursorwidth}
\newlength{\atlasgraphheight}
\newlength{\atlasprecursorheight}
\newlength{\atlasrowheight}
\newlength{\atlasrankwidth}
\newlength{\atlassmilesskip}
% Every column is pinned to a stated width, so the alignment cannot be widened past \textwidth by a
% heading or an index label. These must sum with the three gutters and the divider in the column
% specification: 5mm + 1.8mm + 0.4549 + 2.8mm + 0.5pt + 2.8mm + 0.4549 = \textwidth.
\setlength{\atlasrankwidth}{5mm}
\setlength{\atlasgraphwidth}{0.4549\textwidth}
\setlength{\atlasprecursorwidth}{0.4549\textwidth}
\setlength{\atlasgraphheight}{2.45cm}
\setlength{\atlasprecursorheight}{3.80cm}
% Rows are struck to a common height that exceeds both panels: the surplus is the air between
% structures, and holding it inside the row keeps the column divider unbroken and keeps the
% multipage break in the same place for every row rather than following a wrapped SMILES string.
\setlength{\atlasrowheight}{4.00cm}
\setlength{\atlassmilesskip}{2.4mm}
\newcommand{\atlassmilesstyle}{\scriptsize\ttfamily\color{atlasink}}
\newcommand{\atlasrank}[1]{%
  \makebox[\atlasrankwidth]{\raisebox{-0.5\height}{\footnotesize\color{atlasink}#1}}}
\newcommand{\atlashead}[2]{\makebox[#1]{\footnotesize\scshape #2}}

\title{FORGE: Reaction-Guided Generative Design of Ionizable Lipids\\\large Working draft for 22 assembly families}

\author{Anonymous authors\\Paper under double-blind review}

\begin{document}

\maketitle
\raggedbottom

\begin{abstract}
FORGE generates complete ionizable-lipid graphs conditioned on reaction programs and assesses
their implied assembly by reverse decomposition and exact forward replay. This working draft
extends the evaluation framework from three assembly chemistries to 22 qualified families.
The planned evidence separates exact final assembly (L1), structural plausibility and program
agreement, lipid-reference realism, product and component diversity, held-component generalization,
upstream synthesis evidence (L2) and terminal-material qualification (L3).
\textbf{Final results are pending:} insert family-wise exact-L1 and structural outcomes,
independently trained matched-control contrasts, novelty/diversity, realism and complete-dossier
yields with their denominators and uncertainty. Existing single-checkpoint development diagnostics
and historical three-family results do not establish these new-model claims. This manuscript
addresses computational evidence and makes no new synthesis, formulation or delivery claim.
\end{abstract}

\section{Introduction}

Lipid nanoparticles (LNPs) deliver siRNA, mRNA vaccines and genome-editing reagents
\citep{adams2018patisiran,polack2020safety,baden2021efficacy,gillmore2021crispr}, with ionizable
lipids supporting nucleic-acid complexation, particle assembly and endosomal escape
\citep{cullis2017lipid,hou2021lipid}. Discovery of these lipids commonly relies on combinatorial
synthesis and screening, in which modular precursors are assembled into libraries and the resulting
formulations are evaluated \citep{akinc2008combinatorial,love2010lipidlike,li2024accelerating,xu2024agile}.
This approach uses established reactions to combine known building blocks, which limits the lipids it can generate.

Generative methods incorporate assembly chemistry at different stages of molecular design.
SynNet, RGFN and SynFlowNet retain reaction routes by selecting reactants from a supplied inventory
and applying predefined transformations \citep{gao2022amortized,koziarski2024rgfn,cretu2025synflownet},
a strategy also used by lipid-specific synthesis-DAG and tree-search methods
\citep{ou2024deep,zhao2024generative}. SynCoGen couples precursor design to assembly by jointly
generating building-block graphs, reaction assignments and atomic coordinates
\citep{rekesh2025syncogen}. Whole-molecule models such as DeFoG and GenMol generate structures
without a prescribed precursor inventory \citep{qin2025defog,lee2025genmol}, leaving assembly and
synthesizability to separate assessment \citep{gao2020synthesizability}.

FORGE learns a conditional distribution over complete lipid graphs from reaction programs that
specify roles, reaction-core positions, transformation order and coarse architecture. It receives
no precursor identifiers, stored component graphs or fragment tokens. Reverse decomposition and
exact forward replay verify each generated assembly. A bounded recursive search then assesses
whether the implied precursors have preparation routes supported by reaction and supplier evidence,
returning a computational dossier or an explicit unresolved outcome.

This extension evaluates three questions:
\begin{enumerate}
\item Can one program-conditioned whole-graph model retain exact assembly across all 22 qualified
families under matched held-out requests and independent training seeds?
\item Do the generated structures satisfy independently assessed chemical and architectural checks
while retaining component diversity, novelty and lipid-reference distribution coverage?
\item Which generated precursors and products have complete documentary L2/L3 evidence, and where
does the bounded route assessor remain unresolved or censored?
\end{enumerate}
The results for these questions remain pending. Historical three-family computational results are
retained in explicitly labeled appendices. Final candidate construction may use declared TRAIN-derived
local motif priors; their contribution and cost must be separated from the learned generator and
from independent assessment.

\section{Preliminaries and problem formulation}
\label{sec:preliminaries}

\subsection{Ionizable lipids and molecular architecture}
\label{sec:lipid-background}
Ionizable lipids combine protonatable amine-containing headgroups, linkers and hydrophobic tails
\citep{semple2010rational,jayaraman2012maximizing}. Protonation supports nucleic-acid association
during acidic formulation and membrane interactions involved in endosomal release
\citep{xu2024agile}. LNPs commonly combine an ionizable lipid with helper lipids, cholesterol and
PEGylated lipids. Lipid structure, stereochemistry and formulation composition affect delivery
\citep{han2021ionizable,hatit2023nanoparticle,cheng2020selective}; molecular constitution and
formulated-particle performance are therefore distinct evaluation levels.

\subsection{Reaction families and modular assembly}
\label{sec:reaction-background}
The historical study evaluated three modular assembly patterns. Ugi 3-CR combines an amine, aldehyde and
isocyanide into an $\alpha$-amino amide \citep{xu2024agile}. Repeated aza-Michael addition attaches
acrylate-bearing components to amine sites \citep{li2023combinatorial}; repeated reductive amination
couples aldehyde-bearing components to amines through condensation and reduction
\citep{xue2024barcoding}. The Ugi variant has no carboxylic-acid reactant component and uses an
acidic phosphorus catalyst; its ester functionality comes from the aldehyde precursor.
Repeated assemblies require both transformation order and participating sites. Precursor-origin
boundaries can span functional headgroup, linker and tail regions. These distinctions motivate
role-assigned reaction programs (Appendix~\ref{app:chemistry}, Table~\ref{tab:reaction-family-background}).

\subsection{Molecular graphs and reaction programs}
\label{sec:mathematical-objects}
A lipid constitution is a nonempty, connected, simple molecular graph
$G=(V,E,\ell_V,\ell_E)$, with atom labels $\ell_V:V\to\Sigma_{\rm atom}$ and bond labels
$\ell_E:E\to\Sigma_{\rm bond}$ in finite vocabularies. Bond order is an edge label.
We compare graphs up to label-preserving isomorphism after removing atom-map and stereochemical
labels, retaining formal charges and bond orders. We write $G$ for a graph or its constitutional
isomorphism class when the distinction is immaterial. The bounded graph space
$\mathfrak G_{N,K_{\rm cl}}$ contains these classes with $1\leq|V|\leq N$ and cycle rank
$\beta(G)=|E|-|V|+1\leq K_{\rm cl}$. This is a structural support definition; molecular validity
additionally requires valence and sanitization checks. The historical model used $N_{\max}=194$
and $K_{\rm cl}=3$. The 22-family development path supports up to 254 heavy atoms and 12 closure
edges, including the qualified Br extension; Table~\ref{tab:22-training} will pin the final support
contract. The formal statements use symbolic $N_{\max}$ and $K_{\rm cl}$.

A reaction program specifies the assembly constraints,
\begin{equation}
    P=(c,d,\mathbf r,\mathbf k,M).
    \label{eq:reaction-program}
\end{equation}
Here $c$ identifies a registered assembly chemistry with its admitted transformation sequences,
$d$ is the requested depth, $\mathbf r$ assigns precursor roles to
atom positions, and $\mathbf k$ marks reaction-core positions. The four morphology fields in $M$
record each role's exterior atom count, junction budget, closure count and core-attachment count.
These quantities constrain coarse architecture without specifying precursor identities. Precursor
origin and core membership are separate: an atom can belong to a precursor and participate in the
core; a template-introduced atom belongs to the core without belonging to a precursor.
Appendix~\ref{app:program-encoding} defines the morphology counts.

\subsection{Exact assembly, novelty and route evidence}
For registered chemistry $c$, a precursor input $\mathbf C\in\mathfrak C_c$ is an ordered tuple
of constitutions, roles and transformation steps. Its registered forward products within the
bounded support form $\mathcal F_c(\mathbf C)\subseteq\mathfrak G_{N,K_{\rm cl}}$.
Reverse search returns a finite set $\mathcal D_c(G)$ of precursor inputs and assembly traces
$(\mathbf C,\tau)$. Exact forward replay retains
\begin{equation}
    \widehat{\mathcal E}_c(G)=
    \left\{(\mathbf C,\tau)\in\mathcal D_c(G):G\in\mathcal F_c(\mathbf C)\right\}.
    \label{eq:main-exact-traces}
\end{equation}
Membership uses constitutional identity. A valid graph passes verified exact level-one assembly
(L1) when $\widehat{\mathcal E}_c(G)\neq\varnothing$. Bounded reverse search can miss assemblies,
and recovered traces need not match every sampled depth, role-position or morphology field in $P$.
Assembly existence and conditioning adherence are distinct checks, as are decomposition coverage
and replay precision (Appendix~\ref{app:verification}).

A recovered precursor is component-novel when its constitution is absent from its role's catalogue.
Strong product-level catalogue escape requires \emph{every} exact precursor input to contain an
outside component. We measure recovered component novelty, which does not establish that universal
claim (Appendix~\ref{app:catalogue-reachability}). A complete computational synthesis dossier also
requires evidence-qualified precursor preparation (L2) and terminal-material availability (L3).
Incomplete assessments remain unresolved or search-censored; these checks do not establish
experimental synthesis success or delivery activity.

\section{FORGE}
\label{sec:method}

\noindent\textbf{Implementation scope.} The equations below describe the shared framework.
Historical training and Ugi-specific readout details are retained for traceability; the final
22-family implementation, loss, masks, source priors and repair/selection maps require the contract
in Table~\ref{tab:22-training} and the correspondence audit in Table~\ref{tab:22-theory}.
No historical implementation statement is evidence that a new-model experiment has completed.

\subsection{Overview}
\label{sec:forge-overview}

FORGE learns the distribution $p_\theta(G\mid P)$ of complete lipid molecules $G$ given a
reaction program $P$. The model receives no precursor identifiers, stored component graphs or
fragment tokens. The program specifies how the molecule is assembled without selecting the
precursor structures. FORGE has four main steps (Figure~\ref{fig:forge-overview}):
\begin{enumerate}
    \item \textbf{Specify the reaction program.} Choose assembly chemistry, roles, reaction-core
    positions, transformation order and morphology, leaving precursor identities open.
    \item \textbf{Generate the whole graph.} A conditional discrete flow generates atom, bond and
    topology states while preserving program-fixed coordinates.
    \item \textbf{Verify exact assembly.} Recover role-assigned precursors and require their forward
    transform to reconstruct the generated molecule.
    \item \textbf{Construct precursor routes.} Search for preparation routes to evidence-qualified
    starting materials and abstain when a complete dossier cannot be constructed.
\end{enumerate}

\begin{figure*}[htbp]
\centering
\resizebox{\textwidth}{!}{\input{figures/forge_algorithm_schematic/schematic_body.tex}}
\caption{\textbf{FORGE generation and route assessment.}
A reaction program fixes assembly roles and core coordinates without selecting precursor identities.
FORGE generates the whole graph, verifies exact assembly replay, and recursively traces its implied
precursors to evidence-qualified materials or an explicit abstention. Colours mark precursor origin;
grey denotes the Ugi template-introduced oxygen. Replay and route closure are computational checks,
not experimental synthesis evidence.}
\label{fig:forge-overview}
\end{figure*}

\subsection{Conditioning on reaction programs}
\label{sec:program-conditioning}
An admissible complete program is sampled from the training-fold prior,
\begin{equation}
    p_\theta(G,P\mid c)=p_{\mathrm{prog}}(P\mid c)\,p_\theta(G\mid P).
\end{equation}
Learned embeddings combine chemistry, depth, precursor role, core status and role-local morphology
at each graph position. The adapter determines active positions and fixed or variable coordinate
masks. Appendix~\ref{app:program-encoding} specifies the embeddings and morphology counts.

\subsection{Graph representation and discrete flow}
\label{sec:graph-flow}
FORGE represents each molecule as a spanning tree with additional bonds for ring closures.
A deterministic encoder maps the molecular graph $G$ and its selected atom-mapped
assembly annotation $\alpha$ to
\begin{equation}
    X_\star := \operatorname{Enc}_P(G,\alpha)
    = (\mathbf a,\boldsymbol\pi,\mathbf b,\mathbf e^L,\mathbf e^R,\mathbf z),
\end{equation}
where $\mathbf a$ records atom states, $\boldsymbol\pi$ identifies each atom's parent in the tree,
and $\mathbf b$ records the tree bonds. The arrays $\mathbf e^L$ and $\mathbf e^R$ identify the
endpoints of additional bonds, whose states are stored in $\mathbf z$.

The padded state space $\overline{\mathfrak X}_{N,K_{\rm cl}}$ is finite. Its valid encodings form
$\mathfrak X_{N,K_{\rm cl}}$: each contains a connected tree and at most $K_{\rm cl}$ distinct
closure edges. The decoder $\operatorname{Dec}$ reconstructs the graph from a valid encoding and
returns $\bot_{\rm dec}$ for a malformed state. Every graph within the specified atom, bond, size
and closure limits has an encoding (Theorem~\ref{thm:bounded-representation}).

During sampling, we also restrict the allowed parents and number of children for each atom.
These constraints limit which valid encodings the sampler can use.

For a given $P$, let $\mathcal I(P)$ index the categorical coordinates, each with finite alphabet
$\mathcal X_i(P)$. Pointer alphabets contain only positions permitted by the layout.
The adapter partitions coordinates into fixed positions $\mathcal F(P)$ with values $x_{P,i}$ and
variable positions $\mathcal V(P)$. The set $\Omega_P$ contains raw states within the allowed alphabets that satisfy the fixed
values. It can contain malformed or chemically invalid states. The source law $q_0^P$ fixes
$\mathcal F(P)$ and uses independent, full-support marginals $q_{0,i}^P$ on variable alphabets.
For a clean endpoint $x_1\in\Omega_P$, the conditional corruption law is
\begin{equation}
    q_{t\mid1}^P(x\mid x_1)
    =\mathbf 1[x_{\mathcal F(P)}=x_{P,\mathcal F(P)}]
    \prod_{i\in\mathcal V(P)}
    \left[(1-t)q_{0,i}^P(x_i)+t\mathbf 1\{x_i=x_{1,i}\}\right],
    \quad 0\leq t\leq1.
    \label{eq:conditional-path}
\end{equation}
Each coordinate interpolates between its source and clean value. Their conditional joint law is a
product, so it generally differs from a linear mixture of the joint source and endpoint laws.

The clean-coordinate predictor $d_{\theta,i}(a\mid x,t,P)$ assigns a probability to each
$a\in\mathcal X_i(P)$. A continuous-time Markov chain (CTMC) uses these predictions to average
conditional one-coordinate rates,
\begin{equation}
 R_{t,i}^\theta(u,v\mid x,P)
 =\sum_{a\in\mathcal X_i(P)}d_{\theta,i}(a\mid x,t,P)
 R^*_{t,i}(u,v\mid a,P),\qquad u\neq v,\quad t<1.
 \label{eq:posterior-rate}
\end{equation}
The conditional rate $R^*_{t,i}$ realizes the corresponding factor in
Equation~\ref{eq:conditional-path}; Appendix~\ref{app:flow-proof} defines it and proves the path
identity. The joint generator permits one coordinate to change at an infinitesimal time.
Coordinate posterior marginals therefore suffice to specify its rates, even when endpoint
coordinates are statistically dependent.

\subsection{Program-conditioned Transformer and training}
\label{sec:model-training}

A graph Transformer predicts clean coordinates from $(X_t,t,P)$. Each layer combines graph
self-attention with parent, child and closure relation biases and cross-attention to program,
depth, role and core tokens. Softly program-routed residual adapters add family-specific parameters.
Qualified atom-mapped traces supervise role and core annotations; unresolved source-trace ambiguity
excludes a record from semantic training.

The historical source-balanced training law $\mu_{\rm bal}$ covers annotated triples $(G,\alpha,P)$
with equal mass across the three historical chemistries, not across individual morphology-specific
programs. The recorded historical model combines masked coordinate cross-entropy, role-balanced chemistry, role/core supervision,
and repeat and topology consistency terms. That implementation combined chemistry-task gradients with PCGrad
\citep{yu2020gradient}. Appendix~\ref{app:training} specifies the full objective and training-time
distribution.

To isolate the statistical property of coordinate denoising, define the reference population loss
\begin{equation}
 \mathcal L_{\rm ref}(\theta)
 =\mathbb E_{\substack{(G,\alpha,P)\sim\mu_{\rm bal},\ t\sim\mathcal U(0,1)\\
 X_t\sim q_{t\mid1}^P(\cdot\mid X_\star)}}
 \left[\sum_{i\in\mathcal V(P)}w_i(P)
 \bigl[-\log d_{\theta,i}(X_{\star,i}\mid X_t,t,P)\bigr]\right],
 \qquad w_i(P)>0.
 \label{eq:primary-loss}
\end{equation}
The weights depend only on the program; they can normalize coordinate counts without depending
on the unknown clean category. This reference loss supports the population analysis. It is not
the complete production training objective.

The reference objective has an ideal population interpretation. For a program with positive training
mass, let $q_1^P$ be the law of $\operatorname{Enc}_P(G,\alpha)$ under $\mu_{\rm bal}(\cdot\mid P)$,
and let $\mu_G^P$ be the corresponding graph law.
\begin{theorem}[Ideal conditional generation]
\label{thm:ideal-terminal-law}
Fix such a program $P$, with registered chemistry $c$, finite coordinate alphabets and a
full-support source on each variable alphabet. Assume every training endpoint belongs to $\Omega_P$, decodes to its source
graph, and passes molecular validity and verified exact L1. Assume an unrestricted measurable
coordinate predictor attains a population minimum of Equation~\ref{eq:primary-loss}. Starting from $q_0^P$, exactly simulate the local CTMC
of Equation~\ref{eq:posterior-rate} on $0\leq t<1$. Then its marginal law converges to $q_1^P$ as
$t\uparrow1$, and
\begin{equation}
 \operatorname{Dec}_{\#}q_1^P=\mu_G^P,\qquad
 \Pr_{G\sim\mu_G^P}\!\left[\mathrm{Valid}(G)=1,
 \widehat{\mathrm{L1}}(G,c)=1\right]=1.
\end{equation}
Here $\operatorname{Dec}_{\#}$ denotes the distribution obtained by decoding.
\end{theorem}
Appendix~\ref{app:flow-proof} proves this result by recovering the clean-coordinate posteriors and
verifying the CTMC forward equation. The theorem concerns the reference population objective and the
limit of exact continuous-time sampling. It establishes neither generalization beyond $q_1^P$ nor
correctness of finite training, PCGrad, capped Euler updates or constrained terminal readout.

\subsection{Sampling and exact assembly verification}
\label{sec:sampling-verification}
The numerical sampler uses states $Z_k$, distinct from the ideal CTMC $X_t$. Starting from
$Z_0\sim q_0^P$, it applies
\begin{equation}
 Z_{k+1}=\Pi_P\!\left(\Phi_k^\theta(Z_k,\xi_k)\right),
 \qquad k=0,\ldots,L_{\rm samp}-1,
 \label{eq:numerical-sampling}
\end{equation}
where $\Phi_k^\theta$ is a capped independent-coordinate Euler update, $\xi_k$ contains its random
draws and $\Pi_P$ restores program-fixed coordinates. Fixed coordinates have zero rate and remain
constant (Proposition~\ref{prop:fixed-invariance}); chemical consistency requires verification.
A fresh prediction at $t=1$ drives the constrained argmax readout
$Z_{\rm out}=H_P^\theta(Z_{L_{\rm samp}})$, including Ugi core-saturation rules, or a typed failure.
We then decode, sanitize and apply Equation~\ref{eq:main-exact-traces} to each valid graph.
Molecular validation, decomposition, exact replay and root ambiguity are recorded separately.
Appendix~\ref{app:sampling} specifies the numerical operations and their distinction from the ideal
CTMC.

\subsection{Recursive route assessment}
\label{sec:route-assessment}
For a molecule with exactly one verified precursor decomposition $(\mathbf C_T,\tau_T)$,
we search for routes to prepare its precursors within fixed depth and query limits:
\begin{equation}
    Y=\mathsf{Assess}_{\mathfrak B}(G,P,\mathbf C_T,\tau_T)
    \in\mathfrak D\cup\mathfrak F_{\mathfrak B}.
\end{equation}
Here $\mathfrak B$ is the frozen route contract, $\mathfrak D$ the dossier space and
$\mathfrak F_{\mathfrak B}$ its finite typed outcome set. Ambiguous roots cause abstention.
A complete dossier requires molecular validity, exact L1, admitted and verified upstream reactions,
and available terminal materials. Other outcomes are unresolved or search-censored. Assessment
operates on the completed graph without modifying generation. Appendix~\ref{app:route-assessment}
and Algorithm~\ref{alg:forge-inference} specify admission and inference; Appendix~\ref{app:theory}
states the formal guarantees and their assumptions.

\input{sections/results_22.tex}

\section{Discussion}
\ResultSlot{R03--R10}{State only the conclusions supported by the final 22-family results. Discuss
the weakest families, negative controls, remaining structural defects, generalization limits and
incomplete route evidence before making a breadth or quality claim.}
The intended contribution is computational whole-molecule design with explicit reaction and
synthesis-evidence accounting. High exact-L1 yield alone does not establish realistic lipid
architecture, upstream route completeness, synthesis success or delivery activity. TRAIN-motif
repairs require independent assessment, and gains obtained with additional proposals must be
reported with their cost. Historical three-family findings motivate the extension but cannot
supply its effect sizes or uncertainty. The final discussion will distinguish measured/computed
outcomes, supported interpretation and unresolved hypotheses.

% Main scientific content ends here; ICLR excludes the following disclosure statements.
\clearpage

\subsection*{AI use statement}
Generative AI tools assisted with conceptual and methodological development, software implementation
and debugging, experimental design and diagnostic analysis, interpretation of computational results,
literature search and summarization, and drafting and editing the manuscript. The authors are
responsible for the final code, reported results, citations and conclusions, including material
produced with AI assistance.

\subsection*{Ethics statement}
This 22-family extension is computational. Chemical validity, exact assembly, documentary route
evidence and predicted properties do not establish biological safety or efficacy. No procurement,
synthesis, formulation, animal experiment or candidate lock is initiated by this result scaffold.
Author-supplied experimental material in the copied source remains archived and is not evidence
for the new 22-family model.

\subsection*{Reproducibility statement}
The Methods and Appendix specify the molecular representation, conditional discrete flow, numerical
sampler, exact assembly verification, route assessment, benchmark protocols and evaluation metrics.
Appendix~\ref{app:repro} summarizes randomization, historical implementation provenance and the
access limitations of the standalone manuscript archive. Historical seed-level results appear in Appendix~\ref{app:results}; new-model seed rows remain
pending in Appendix~\ref{app:22-results}. Assumptions and proofs appear in Appendix~\ref{app:theory}.

\bibliographystyle{iclr2027_conference}
\bibliography{forge,references}

\clearpage
\appendix

\input{sections/appendix_22.tex}
\clearpage

\section{Shared framework and historical implementation details}
\label{app:forge}

This appendix specifies program and graph encoding, model training, sampling, verification and
route assessment. Algorithm~\ref{alg:forge-inference} combines the inference steps.

\subsection{Reaction-program encoding}
\label{app:program-encoding}

The program tuple in Equation~\ref{eq:reaction-program} uses the ordered atom-position set
$\mathcal I_{\rm atom}(P)=\{1,\ldots,n(P)\}$. The arrays $\mathbf r$ and $\mathbf k$ assign
precursor roles and core positions to these atoms. This atom-position index differs from
$\mathcal I(P)$, which indexes all six families of categorical coordinates.

Morphology is computed from the selected atom-mapped trace and its sparse serialization.
For role $r$, let $V_r^{\rm ext}$ contain its exterior atoms, excluding reaction-core atoms, and
let $o_r(v)$ count the same-role exterior children of $v$ in the serialized tree. The four fields are
\begin{equation}
 M_r=(n_r,j_r,\gamma_r,a_r),\qquad
 n_r=|V_r^{\rm ext}|,\qquad
 j_r=\sum_{v\in V_r^{\rm ext}}\max\{o_r(v)-1,0\}.
\end{equation}
Here $\gamma_r$ counts closure edges with both endpoints in $V_r^{\rm ext}$, and $a_r$ counts
same-role bonds between exterior and core atoms. When the serialized tree restricts to a spanning
forest of the exterior subgraph, $\gamma_r$ equals that subgraph's cycle rank. The junction budget
is a tree-branching count; it does not count degree contributed by closure bonds. Consequently these
fields depend on the selected annotation and serialization, not on the unannotated graph alone.
Fixed-coordinate equality alone does not establish agreement with every morphology field.
Positions without precursor origin use the designated absent morphology state. Integer morphology
values are encoded with an offset of one, reserving zero for that absent state.

The training-fold program prior is defined in Section~\ref{sec:program-conditioning}.
At graph position $i$, the program embedding is
\begin{equation}
    h_i^P=\operatorname{LN}\!\left(
        e_c(c)+e_d(d)+e_r(r_i)+e_k(k_i)
        +\sum_{u=1}^{4}e_{M,u}(M_{r_i,u})
    \right).
\end{equation}
Here, $e_c,e_d,e_r,e_k$ and $e_{M,u}$ are learned embeddings in a shared hidden space, and
$\operatorname{LN}$ denotes layer normalization. The embedding encodes precursor role, reaction-core
status and role-local morphology. The morphology fields determine active positions and fixed or
variable masks. Sampling from the prior specifies these constraints; generation assigns the atom,
bond and topology states within each role.

A chemistry-specific adapter determines the precursor roles, reaction core, reverse decomposition
and forward transform associated with $c$. Its sampling layout determines active positions and
fixed or variable coordinates; the embeddings condition neural prediction on the remaining
program fields.

\subsection{Graph serialization and source annotations}
\label{app:graph-encoding}

The historical spanning-tree representation used at most three residual closure edges and
194-atom support without an $O(N_{\max}^2)$ edge state. These historical numerical limits do not
replace the larger 22-family support contract in Table~\ref{tab:22-training}. The discrete flow generates the atom, bond
and topology states in this representation.
Let $\alpha=(\mathbf C,\tau,\mathbf r,\mathbf k)$ contain the selected exact precursor tuple,
registered assembly trace, atom-wise precursor-role assignment and reaction-core annotation. The deterministic
program-aware encoder
\begin{equation}
    \operatorname{Enc}_P:
    \{(G,\alpha):G\in\mathfrak G_{N_{\max},K_{\rm cl}}
      \text{ and }\alpha\text{ is admitted for }P\}
    \longrightarrow \mathfrak X_{N_{\max},K_{\rm cl}}
\end{equation}
constructs the categorical graph state
\begin{equation}
    X_\star:=\operatorname{Enc}_P(G,\alpha)=
    \left(
        \mathbf{a},
        \boldsymbol{\pi},
        \mathbf{b},
        \mathbf{e}^{L},
        \mathbf{e}^{R},
        \mathbf{z}
    \right)
\end{equation}
where $\mathbf{a}$ contains atom states, $\boldsymbol{\pi}$
parent pointers, $\mathbf{b}$ spanning-tree bond states, $\mathbf{e}^{L}$ and $\mathbf{e}^{R}$
closure-endpoint pointers, and $\mathbf{z}$ closure-bond states.

The encoder first partitions the atoms into connected precursor-origin components under $\alpha$.
It selects as root the reaction-core atom with the smallest tie-broken canonical RDKit rank, and
orders the origin-component graph by breadth-first traversal from that root component, breaking
ties by role index, component canonical ranks and original atom index. Within each component it uses
a deterministic depth-first preorder from the selected attachment atom, with neighbors ordered by
the same canonical ranks. The traversed bonds define the spanning tree. Every residual bond is
represented once as an ordered endpoint pair with its bond label, and residual triples are sorted
lexicographically. Role, core, fixed-coordinate and morphology arrays follow the
same atom order. We excluded training records with more than one qualified exact source trace
from semantic training and collapsed duplicate descriptions of an admitted trace to its canonical
record. Thus
$\operatorname{Dec}(\operatorname{Enc}_P(G,\alpha))\cong G$ for every admitted record.

The raw padded categorical product space $\overline{\mathfrak X}_{N_{\max},K_{\rm cl}}$ also
contains malformed states. Its well-formed subset $\mathfrak X_{N_{\max},K_{\rm cl}}$ enforces an
active node prefix, earlier-parent pointers, distinct valid closure endpoints, no duplicate edges,
and null values in unused closure slots. The total decoder, including its failure value, is typed as
\begin{equation}
    \operatorname{Dec}:\overline{\mathfrak X}_{N_{\max},K_{\rm cl}}
    \longrightarrow
    \mathfrak G_{N_{\max},K_{\rm cl}}\cup\{\bot_{\rm dec}\}.
\end{equation}
For the historical three-family study, $N_{\max}=194$ and $K_{\rm cl}=3$.
These bounds specify architectural support. Reaction consistency, nonzero probability under a
trained model and synthesis-program completeness require separate conditions.

The structural decoder specifies graph reconstruction from an already chosen state.
The production terminal readout additionally chooses that state from neural predictions
(Appendix~\ref{app:sampling}). These are separate operations.

The chemistry adapter marks any atom, pointer, or bond states fixed by $P$.
FORGE holds these states constant while corrupting the remaining states along the linear probability
path and predicting their clean values from $(X_t,t,P)$. Let $\Omega_P$ be the subset of the raw state
space satisfying the allowed coordinate alphabets and fixed values, and define the admissible program set
\begin{equation}
    \mathcal P_{\rm adm}
    =\left\{P:\mathfrak X_{N_{\max},K_{\rm cl}}\cap\Omega_P\neq\varnothing\right\}.
\end{equation}
The morphology prior is normalized only on $\mathcal P_{\rm adm}$. Each chemistry adapter is required
to satisfy the program-compatible encoding property: every admitted annotated record obeys
$\operatorname{Enc}_P(G,\alpha)\in\Omega_P$. During corpus construction, we checked this property and exact constitutional decode round trip
for every record before split assignment. Records failing either check were excluded. The checks
covered all cached training, calibration and held-out records; they do not establish the property
for unregistered chemistry.

\subsection{Transformer architecture and training objectives}
\label{app:training}

The denoiser is a graph Transformer over the complete sparse state. Each layer applies graph
self-attention with relation biases for parent, child and closure edges, then cross-attends to memory
tokens representing program identity, program depth, precursor role and reaction-core position.
Cross-attention is repeated in every layer so that reaction semantics can influence both local state
updates and long-range coordination among precursor-derived regions. Three softly program-routed
residual adapters add family-specific parameters to the shared whole-graph backbone.

The reported final three-family model used a stratified minibatch with 21 records per chemistry
and source-weighted sampling within each stratum. Two such microbatches formed an optimizer update.
Training times were $t=\min\{0.98,\max\{0.02,U\}\}$ for $U\sim\mathcal U(0,1)$, giving point
masses of $0.02$ at both endpoints and unit density on $(0.02,0.98)$. Thus training did not cover
all times required by the ideal reference objective in Equation~\ref{eq:primary-loss}.

For chemistry $c$ and state family
$q\in\mathscr J=\{\mathbf a,\boldsymbol\pi,\mathbf b,\mathbf e^L,\mathbf e^R,\mathbf z\}$,
let $\mathcal S_{c,q}$ be the selected variable coordinates across that chemistry's minibatch.
With coordinate loss $\ell_i=-\log d_{\theta,i}(X_{\star,i}\mid X_t,t,P)$, the pooled reduction is
\begin{equation}
 L_{c,q}^{\rm pool}=\frac{\sum_{i\in\mathcal S_{c,q}}\ell_i}{\max\{1,|\mathcal S_{c,q}|\}}.
\end{equation}
Coordinates here include their example index. Chemistry states (atoms, parent bonds and closure
bonds) additionally use equal mass over roles present in the selected coordinates. If
$\mathcal S_{c,q,r}$ contains coordinates assigned to role $r$ and
$\mathcal R_{c,q}=\{r:|\mathcal S_{c,q,r}|>0\}$, this reduction is
\begin{equation}
 L_{c,q}^{\rm role}=\frac{1}{|\mathcal R_{c,q}|}
 \sum_{r\in\mathcal R_{c,q}}\frac{\sum_{i\in\mathcal S_{c,q,r}}\ell_i}{|\mathcal S_{c,q,r}|},
 \label{eq:role-balanced-loss}
\end{equation}
with zero loss if no role is present. Atom and parent-bond groups use their position's role;
closure-bond groups use the role at the clean left endpoint. Pointer losses use the pooled reduction
and mask logits to the permitted layout positions. The flow term $L_{c,\rm flow}$ sums these six
reduced losses, with unit weights.

For each chemistry, the recorded task loss was
\begin{equation}
\begin{aligned}
 L_c={}&L_{c,\rm flow}
 +0.25L_{c,\rm role}+0.25L_{c,\rm core}+0.25L_{c,\rm repeat}\\
 &+L_{c,\rm child}+0.5L_{c,\rm junction}+L_{c,\rm chem\mid top}.
\end{aligned}
\label{eq:production-task-loss}
\end{equation}
The role and core terms are cross-entropies averaged equally over present target states.
The repeat term averages squared differences between predicted atom distributions, and separately
parent-bond distributions, at aligned exterior positions in repeated component instances; it sums
the two means. The child term is state-balanced cross-entropy for same-component exterior child
counts in $\{0,1,2,3\}$. The junction term applies smooth-L1 loss to predicted and target
role-local junction totals, each divided by the number of exterior atoms in that role, then averages
over eligible examples within a role and equally across present roles. Predicted junction totals
sum the expected child-count contribution $\max\{o-1,0\}$ across exterior atoms.
The topology-conditioned chemistry term repeats the role-balanced atom and bond losses in a second
forward pass with clean parent and closure endpoints and corrupted atom and bond states.
These definitions describe the reported final arm; the earlier architecture study used its own
recorded ablation settings.

PCGrad combines the task gradients in integer chemistry-state order. For
$g_c=\nabla_\theta L_c$, it initializes $\widetilde g_c=g_c$ and, for each other task $h$ in that
order, applies
\begin{equation}
 \widetilde g_c\leftarrow\widetilde g_c-
 \mathbf1\{\langle\widetilde g_c,g_h\rangle<0\}
 \frac{\langle\widetilde g_c,g_h\rangle}{\max\{\|g_h\|_2^2,10^{-12}\}}\,g_h
\end{equation}
on parameter coordinates available to both tasks. Gradients are not norm-normalized, and the
optimizer receives the arithmetic mean of the projected task gradients. This update is not generally
the gradient of a scalar sum of task losses. No population optimality claim is made for the full
procedure. In particular, target-dependent group reductions and the restricted training-time range
require analysis beyond Proposition~\ref{prop:masked-bayes}.

We tested layerwise cross-attention, each auxiliary loss, routed adapters and gradient-conflict
control through separate interventions (Appendix~\ref{app:ablation-results}).

\subsection{Numerical sampling and terminal decoding}
\label{app:sampling}

The finite sampler evaluates the denoiser at $t_k=k/L_{\rm samp}$, for
$k=0,\ldots,L_{\rm samp}-1$, and uses $\Delta t=1/L_{\rm samp}$.
Conditional on $Z_k=x$, it draws one clean category
$A_{k,i}\sim d_{\theta,i}(\cdot\mid x,t_k,P)$ independently for each variable coordinate.
Write $s_{t,i}^P(u\mid a)=(1-t)q_{0,i}^P(u)+t\mathbf1\{u=a\}$ for its conditional path factor,
$\dot s_i^P(u\mid a)=\mathbf1\{u=a\}-q_{0,i}^P(u)$ for the derivative, and
$[z]_+=\max\{z,0\}$. The code uses the conditional off-diagonal rate
\begin{equation}
 R_{t,i}^{*,\varepsilon}(u,v\mid a,P)
 =\frac{[\dot s_i^P(v\mid a)-\dot s_i^P(u\mid a)]_+}
 {|\mathcal X_i(P)|\max\{s_{t,i}^P(u\mid a),\varepsilon\}},
 \quad u\neq v,\qquad \varepsilon=10^{-8}.
 \label{eq:clamped-rate}
\end{equation}
Full source support is required only on the allowed coordinate alphabet, including the masked
candidate set for a pointer. Thus that alphabet is the support for every $t_k<1$.
The denominator clamp in Equation~\ref{eq:clamped-rate} is a numerical safeguard absent from the
ideal rate in Appendix~\ref{app:flow-proof}.

For the sampled endpoint $a=A_{k,i}$, let
$h_{k,i}=\Delta t\sum_{v\neq x_i}R_{t_k,i}^{*,\varepsilon}(x_i,v\mid a,P)$ and set
$\kappa_{k,i}=\min\{1,0.999/h_{k,i}\}$ when $h_{k,i}>0$, with $\kappa_{k,i}=1$ otherwise.
The implemented coordinate kernel is
\begin{equation}
 \widehat K_{k,i}(v\mid x,a)=
 \begin{cases}
 \kappa_{k,i}\Delta t\,R_{t_k,i}^{*,\varepsilon}(x_i,v\mid a,P),&v\neq x_i,\\
 1-\kappa_{k,i}h_{k,i},&v=x_i.
 \end{cases}
 \label{eq:capped-kernel}
\end{equation}
Its off-diagonal masses sum to at most $0.999$; the diagonal is the remaining mass. Consequently
all entries are nonnegative and each row sums to one. Coordinates are drawn independently from
these kernels, followed by the projection $\Pi_P$.

Without the clamp or cap, and when every conditional Euler row is nonnegative, averaging over
$A_{k,i}$ gives the Euler kernel built from Equation~\ref{eq:posterior-rate} exactly. One sampled
endpoint per coordinate therefore does not by itself bias that uncapped one-step law. The nonlinear
cap generally changes the averaged kernel. Independent coordinate updates also permit simultaneous
finite-step changes, so their product is a numerical approximation to the local CTMC, not its exact
transition kernel. No discretization error bound is asserted near the terminal singularity.

After all steps, the model evaluates $(Z_{L_{\rm samp}},1,P)$ once more. The production terminal
readout $H_P^\theta$ selects coordinate values using constrained argmax, restores fixed states and
returns either a raw state or a typed terminal failure, denoted collectively by $\bot_{\rm term}$.
For Ugi it enforces one hydrogen on the aldehyde-derived $\alpha$ carbon and one on the
isocyanide-derived amide nitrogen. The selected state is then reconstructed and checked for
connectedness, valence and sanitization. A readout failure retains its detailed reason in the ledger;
for graph-law notation we extend $\operatorname{Dec}(\bot_{\rm term})=\bot_{\rm dec}$.
If $q_{\theta,\rm out}^P$ is the law after terminal readout, including failure, then
\begin{equation}
 p_\theta(\cdot\mid P)=\operatorname{Dec}_{\#}q_{\theta,\rm out}^P,
 \qquad
 p_\theta(G\mid P)=\sum_{z:\operatorname{Dec}(z)=G}q_{\theta,\rm out}^P(z).
 \label{eq:generated-law}
\end{equation}
This law retains structural failures and chemically invalid graphs before the validity and L1
filters; it is not conditioned on acceptance.
The ideal theorem does not include the learned $t=1$ readout, which is unconstrained by a loss
integrated over $t\in(0,1)$. We compared core-saturation and valence-only readouts under the same
trained weights and verifier, then separately compared program-role and global source marginals
(Appendix~\ref{app:ablation-results}).

\subsection{Exact assembly verification}
\label{app:verification}

For the chemistry $c$ specified by $P$, FORGE decomposes each valid graph into role-assigned
precursors and applies the registered forward transform. Exact replay checks assembly existence
under $c$; it does not test full agreement with the sampled conditioning layout.

\begin{definition}[Exact assembly and recovered witnesses]
\label{def:exact-assembly}
For the registered forward operator $\mathcal F_c$ on precursor inputs $\mathfrak C_c$, define
\begin{equation}
 \mathcal E_c^\star(G)=\{\mathbf C\in\mathfrak C_c:G\in\mathcal F_c(\mathbf C)\},\qquad
 \mathrm{L1}^\star(G,c)=\mathbf1\{\mathcal E_c^\star(G)\neq\varnothing\}.
\end{equation}
The bounded reverse procedure returns candidate pairs $(\mathbf C,\tau)$ in $\mathcal D_c(G)$.
Equation~\ref{eq:main-exact-traces} retains the pairs passing exact forward replay, and
\begin{equation}
 \widehat{\mathrm{L1}}(G,c)=\mathbf1\{\widehat{\mathcal E}_c(G)\neq\varnothing\}.
\end{equation}
Both predicates describe assembly existence; molecular validity is checked separately.
\end{definition}
The set $\mathcal E_c^\star(G)$ contains precursor inputs, whereas
$\widehat{\mathcal E}_c(G)$ contains input--trace pairs. Their compatible comparison is
\begin{equation}
 \{\mathbf C:\exists\tau,\ (\mathbf C,\tau)\in\widehat{\mathcal E}_c(G)\}
 \subseteq\mathcal E_c^\star(G),\qquad
 \widehat{\mathrm{L1}}(G,c)\leq\mathrm{L1}^\star(G,c).
\end{equation}
The inclusion follows directly from exact replay. Its converse requires complete reverse search,
which has not been established. Throughout the paper, verified exact-L1 yield counts valid
attempts satisfying the recovered predicate. Candidate traces are compared after the registered
role- and step-aware canonicalization, so duplicate encodings do not create additional witnesses.

Post-hoc filtering changes which completed samples are retained but cannot change the proposal
distribution's exact-L1 acceptance probability.
Under independent sampling, this probability determines the complete distribution of accepted
counts under a fixed attempt budget and the expected cost of obtaining a fixed number of accepted
products (Proposition~\ref{prop:posthoc-budget}).
We distinguish decomposition coverage from forward-replay precision because a candidate disconnection
may fail to reconstruct $G$.
Candidate decomposition coverage is the fraction of valid graphs for which
$\mathcal D_c(G)\neq\varnothing$. For a declared sequence of valid evaluation graphs $G_1,\ldots,G_m$, candidate replay precision is
\begin{equation}
 \mathrm{ReplayPrecision}
 =\frac{\sum_{b=1}^m|\widehat{\mathcal E}_c(G_b)|}
 {\sum_{b=1}^m|\mathcal D_c(G_b)|},
\end{equation}
with N/E when the denominator is zero. Replay of tuples already in
$\widehat{\mathcal E}_c(G)$ has $100\%$ success by construction. This invariant does not estimate
candidate replay precision. We separately define
\begin{equation}
 \mathrm{CandidateAmbig}(G,c)=\mathbf 1\{|\mathcal D_c(G)|>1\},\qquad
 \mathrm{ExactAmbig}(G,c)=\mathbf 1\{|\widehat{\mathcal E}_c(G)|>1\}.
\end{equation}
A distinct exact-L1 product is a deduplicated generated constitution. An unambiguous exact
decomposition satisfies $|\widehat{\mathcal E}_c(G)|=1$. We distinguish product uniqueness from
decomposition uniqueness throughout.

\subsection{Recursive route assessment and dossier admission}
\label{app:route-assessment}

\subsubsection{Assessment interface}

For a valid generated graph $G$ and reaction program $P$ with chemistry $c$, route assessment
returns either a
complete synthesis dossier or a typed abstention.
Let $\mathfrak B$ denote the frozen route contract: depth and query budgets, registered transforms,
evidence-admission rules and a dated supplier snapshot. Let $\mathfrak D$ be the space of candidate
dossiers, each containing a synthesis tree and its reaction and terminal evidence, and let
$\mathfrak F_{\mathfrak B}$ be the disjoint finite set of typed failures.
Generation and assessment are defined by
\begin{equation}
    Y=\mathsf{Assess}_{\mathfrak B}(G,P,\mathbf C_T,\tau_T)
    \in\mathfrak D\cup\mathfrak F_{\mathfrak B}.
\end{equation}
$\mathsf{Assess}_{\mathfrak B}$ is invoked after the verifier returns one admitted root trace
$(\mathbf C_T,\tau_T)\in\widehat{\mathcal E}_c(G)$. It verifies the final assembly, searches for
preparation routes to unavailable precursors and checks terminal supplier evidence. $P$ and the assessor act at different stages: $P$ guides generation, whereas the assessor
constructs a dossier from the completed graph without modifying it.
FORGE accepts $(G,T)$ only when $G$ passes molecular and exact-assembly verification and $T$ traces
every required precursor to an available starting material; otherwise, it abstains.

\subsubsection{Recursive expansion and completion checks}

Completing a synthesis dossier requires evidence for every recovered precursor. Exact L1
verification specifies the precursor structures; availability and preparation routes require
additional evidence. Each precursor must therefore match an available material or have a supported
route from available starting materials.
Let $T$ be a candidate dossier with bipartite synthesis tree $(V_M,V_R,E_T)$ rooted at $G$, where each product molecule
points to its preparation reaction and each reaction points to its required reactants.
The assessor constructs these routes recursively and withholds a complete dossier when any branch
remains unresolved.
A molecule node terminates when a current, dated supplier record lists an exact constitutional match
as available; otherwise, FORGE must expand that node through a supported, forward-verified reaction.
FORGE uses a bounded level-two (L2) planner to propose each expansion.
For each proposal, FORGE records the reaction, reactants, product, and supporting evidence, and
accepts the step only when the forward transform reconstructs the target precursor exactly.
Each accepted reaction adds its reactants as new molecule nodes, and FORGE applies the same
termination-or-expansion rule recursively.
The current assessor admits only an unambiguous verified exact-L1 root. Every dossier therefore
records the sole selected root trace
$(\mathbf C_T,\tau_T)\in\widehat{\mathcal E}_c(G)$; an ambiguous exact set produces a typed
abstention. For each upstream reaction node $r$, let
$\mathrm{Admitted}_{\mathfrak B}(r)=1$ when its reaction and substrate scope satisfy the frozen
evidence policy, and let $\mathrm{Verify}_{\mathfrak B}(r)=1$ when
$\operatorname{target}(r)\in\mathcal F_r(\operatorname{reactants}(r))$.
For each leaf molecule $\ell$, let $\mathrm{Available}_{\mathfrak B}(\ell)=1$ when it satisfies the terminal
supplier-evidence rule.
Let $V_R^{\mathrm{up}}(T)$ denote the upstream preparation-reaction nodes and $L(T)$ the leaf
molecules. $\mathrm{WellFormed}(G,T,P)$ requires that $T$ is a connected, acyclic bipartite tree
rooted at $G$; each internal molecule has exactly one recorded producing reaction; no molecular constitution repeats along an ancestor path; each reaction's
molecule children equal its recorded reactants with their multiplicities; and $T$ has no disconnected, cyclic or
extraneous node. Its distinguished final-assembly record must have product $G$ and precursor
children exactly $\mathbf C_T$ under trace $\tau_T$.
\begin{equation}
\begin{aligned}
    \mathrm{Complete}_{\mathfrak B}(G,T,P)
    ={}&\mathrm{Valid}(G)\,\mathrm{WellFormed}(G,T,P)\\
    &\times\mathbf1\!\left[\widehat{\mathcal E}_c(G)=\{(\mathbf C_T,\tau_T)\}\right]\\
    &\times\prod_{r\in V_R^{\mathrm{up}}(T)}
    \mathrm{Admitted}_{\mathfrak B}(r)\mathrm{Verify}_{\mathfrak B}(r)
    \prod_{\ell\in L(T)}\mathrm{Available}_{\mathfrak B}(\ell).
\end{aligned}
\end{equation}
All factors are binary and empty products equal one. Repeated use of a material is represented by
separate tree occurrences; equality of their constitutions does not merge their nodes.
Under finite search bounds and terminating evidence queries, the assessor terminates and returns a
dossier only when $\mathrm{Complete}_{\mathfrak B}(G,T,P)=1$
(Lemma~\ref{lem:planner-termination} and Invariant~\ref{inv:dossier-soundness}).
An abstention means that FORGE did not complete a dossier under $\mathfrak B$; it does not imply that
no synthesis route exists outside the frozen search space or evidence snapshot.
The guarantee is computational and does not establish experimental synthesis success.

\subsubsection{Recorded evidence and support tiers}

For each accepted pair $(G,T)$, the synthesis dossier contains every molecule and reaction in $T$,
the evidence for each reaction and a dated supplier record for each terminal material.
We classify assembly support relative to the fixed reference space into four tiers:
\textbf{E0 (enumerated):} the generated lipid matches a product that was already constructed
computationally from known components using a supported reaction.
\textbf{E1 (known components):} the generated lipid was not included in the bounded E0 enumeration,
but a supported reaction assembles it entirely from known components.
\textbf{E2 (generated precursor):} a supported reaction assembles the lipid, but at least one
precursor is newly generated and requires a complete L2 route to available starting materials.
\textbf{E3 (new assembly family):} the lipid requires a final assembly reaction outside the
qualified reaction registry.

\subsection{End-to-end inference procedure}
\label{app:inference}

Algorithm~\ref{alg:forge-inference} records every generation attempt, including failures in
molecular validation, exact replay or route completion.

% Ruled caption and procedure; paragraph spacing is internal to the algorithmic list.
\begingroup
\setlength{\parskip}{0pt}
\hrule height 0.8pt\kern 2pt
\refstepcounter{algorithm}
\noindent\textbf{Algorithm~\thealgorithm} FORGE generation and fail-closed synthesis-program construction.
\label{alg:forge-inference}
\par\nobreak\kern 2pt\hrule height 0.4pt\nobreak
\begin{algorithmic}[1]
\Require trained coordinate predictors, chemistry $c$, program prior $p_{\rm prog}(\cdot\mid c)$,
integers $L_{\rm samp}\geq1$ and $B\geq0$, recorded seed $s$, route contract $\mathfrak B$
\Ensure attempt-indexed dossier sequence $\mathcal D_{\rm out}$, exact-L1 ledger $\mathcal L_{\rm out}$,
and typed failure ledger $\mathcal H$
\State initialize the random-number streams from $s$;
$\mathcal D_{\rm out}\gets[\,]$; $\mathcal L_{\rm out}\gets[\,]$; $\mathcal H\gets[\,]$
\For{$b=1,\ldots,B$}
    \State sample and record $P\sim p_{\rm prog}(\cdot\mid c)$
    \State sample $Z_0\sim q_0^P$ and set $Z\gets\Pi_P(Z_0)$
    \For{$k=0,\ldots,L_{\rm samp}-1$}
        \State $t_k\gets k/L_{\rm samp}$; $\Delta t_k\gets1/L_{\rm samp}$
        \State $Z\gets\Pi_P\!\left(\Call{CappedCoordinateEuler}{Z,t_k,\Delta t_k,P}\right)$
    \EndFor
    \State $Z_{\rm out}\gets H_P^\theta(Z)$ \Comment{fresh prediction at $t=1$, constrained readout}
    \If{$Z_{\rm out}$ is a terminal failure}
        \State record $(b,\textsc{TerminalFailure},\text{readout reason})$ in $\mathcal H$; \textbf{continue}
    \EndIf
    \State $G\gets\operatorname{Dec}(Z_{\rm out})$
    \If{$G=\bot_{\rm dec}$ or $\mathrm{Valid}(G)=0$}
        \State record $(b,\textsc{InvalidGraph},\text{decode or validation reason})$ in $\mathcal H$
        \State \textbf{continue}
    \EndIf
    \State $\widehat{\mathcal{E}}_c(G)\gets
    \{(\mathbf C,\tau)\in\mathcal D_c(G):\exists H\in\mathcal F_c(\mathbf C),\ H\cong G\}$
    \If{$\widehat{\mathcal{E}}_c(G)=\varnothing$}
        \State $\mathcal H\gets\mathcal H\mathbin{\|}[(b,\textsc{L1Failure})]$
        \State \textbf{continue}
    \EndIf
    \State $\mathcal L_{\rm out}\gets\mathcal L_{\rm out}\mathbin{\|}
    [(b,G,\widehat{\mathcal E}_c(G))]$
    \If{$|\widehat{\mathcal E}_c(G)|\neq1$}
        \State $\mathcal H\gets\mathcal H\mathbin{\|}[(b,\textsc{AmbiguousL1})]$
        \State \textbf{continue}
    \EndIf
    \State select the sole $(\mathbf C_T,\tau_T)\in\widehat{\mathcal E}_c(G)$
    \State $Y\gets\mathsf{Assess}_{\mathfrak B}(G,P,\mathbf C_T,\tau_T)$
    \If{$Y$ is a dossier $T$ and $\mathrm{Complete}_{\mathfrak B}(G,T,P)=1$}
        \State $\mathcal D_{\rm out}\gets\mathcal D_{\rm out}\mathbin{\|}
        [(b,T)]$
    \ElsIf{$Y\in\mathfrak F_{\mathfrak B}$}
        \State $\mathcal H\gets\mathcal H\mathbin{\|}[(b,Y)]$
    \Else
        \State $\mathcal H\gets\mathcal H\mathbin{\|}[(b,\textsc{IncompleteDossier})]$
    \EndIf
\EndFor
\State \Return $\mathcal D_{\rm out},\mathcal L_{\rm out},\mathcal H$
\end{algorithmic}
\nopagebreak\hrule height 0.8pt
\endgroup

\subsection{Reproducibility and artifact provenance}
\label{app:repro}

The computational results use the historical shared three-family model. Training-seed labels
0, 1 and 2 correspond to random-number seeds 20260825, 20260826 and 20260827, respectively
(Table~\ref{tab:production-seed-counts}). Each comparison retains its declared training split,
attempt budget, checkpoint rule and evaluator. The manuscript provenance records identify the
source artifacts for the reported tables and figures through SHA-256 hashes.

The historical source verification matched the complete seed-0 evaluation-source fingerprint to
a retained code snapshot. No exact match was recovered for the seed-1 and seed-2 evaluation-source
fingerprint among the retained snapshots inspected, and a complete training-source manifest was
not recovered. These limits and the checked per-seed settings are recorded in the accompanying
\nolinkurl{mathematical_review.json}. Current development code does not establish the historical
implementation used for those runs.

The standalone manuscript archive contains the LaTeX source, generated tables, figure assets and
provenance records. It does not bundle the training corpora, model checkpoints, full evaluation
ledgers or a complete historical training-environment and compute manifest. Independent end-to-end
reruns therefore require access to the corresponding historical run artifacts. The copied source
and its author-supplied experimental material remain archived separately; they are not new-model
results. The 22-family result slots and evidence ledger specify the missing final artifacts.

\clearpage

\section{Historical three-family evaluation and additional results}
\label{app:results}

All numerical results in this appendix refer to the historical three-family study. New-model
measurements belong in Appendix~\ref{app:22-results}, with a separate final evidence contract.

The evaluation protocol applies to Section~\ref{sec:results}. Additional analyses quantify
decomposition outcomes, variation across training seeds, component diversity and the bounded guidance
diagnostics. Each comparison includes its seed-level results and metric definitions.

\FloatBarrier
\subsection{Evaluation protocol}
\label{app:evaluation-protocol}

\subsubsection{Data, splits and generation budgets}

We evaluated three assembly programs with different amounts of supervision: Ugi, repeated
aza-Michael addition and repeated reductive amination. The comparisons tested shared generation,
reaction conditioning and recovery of precursors absent from training. We trained three matched shared-model arms: correct program conditioning, no program
information and cyclically incorrect reaction labels. All arms used the same architecture and
optimization budget. We evaluated the prespecified final checkpoint from three independent training
seeds under equal generation-attempt budgets per program.
Our primary metric is verified exact-L1 yield per generation attempt: the number of generated graphs that
pass molecular validity, program-aware decomposition, and exact forward replay divided by the total
number of sampling attempts.
Secondary metrics separate graph quality (validity and connectedness), reaction recovery (verified
coverage, the replay invariant, abstention, and exact ambiguity), and exploration (distinct-product fraction, whole-product and
component novelty, molecular diversity, and effective component count).
We report exact counts, means across seeds, sample standard deviations, individual seed values, and
the minimum and maximum of the three paired seed contrasts. These are descriptive across-seed
summaries, not confidence intervals. Generated molecules are not treated as independent training
replicates. Conditional binomial variation across a fixed model's
\ForgeProseAttemptsPerProgram{} attempts is distinct from
variation induced by retraining.
Training used only the component-disjoint training folds. We used calibration data for diagnostics
and excluded held-out products from training and checkpoint selection. All three shared arms
assigned equal total sampling mass to each program and preserved the frozen within-program source
weights.
We also checked the prespecified numerical acceptance criteria. With three training seeds, these
checks describe the observed runs and do not establish population-level non-inferiority. Zero program-fixed-coordinate violations are an implementation check required by
Proposition~\ref{prop:fixed-invariance}, not learned behavior.

\subsubsection{Comparators and matched controls}

The null and cyclic controls tested different aspects of conditioning. Null generation followed
by exact post-hoc filtering measured performance without program information. The cyclic arm
isolated the reaction-identity label with the remaining program structure retained. A separate
train-only finite-component assembler received the exact training components, supported forward
chemistry, source-weighted role marginals and the same number of attempts. We evaluated reaction
validity and exploration separately because catalogue assembly has zero input-component novelty
by construction. None of these comparisons called an L2 planner, biological oracle or
candidate-selection procedure.

For Ugi, we evaluated a common external-baseline protocol with one
reaction-space comparator, RGFN, and two generic whole-molecule comparators, unconditional DeFoG
and GenMol/SAFE \citep{koziarski2024rgfn,qin2025defog,lee2025genmol}.
The Synthesis-DAG method of \citet{ou2024deep}, SynFlowNet and SynCoGen were discussed as related methods and were
excluded from the empirical comparison \citep{cretu2025synflownet,rekesh2025syncogen}.
At the version checks used to construct the benchmark, the Synthesis-DAG repository lacked the
paper-specific ionizable-lipid implementation, the SynFlowNet action space did not support the
qualified three-reactant Ugi transform, and the SynCoGen repository lacked a usable code license.
We preserved each method's native reaction operations and included numerical results only from
executions using the common training split, attempt budget, seed set and molecular and exact-L1
evaluator. The comparison excludes published scores from other datasets and unsupported
aza-Michael or reductive-amination runs.
The matched internal suite included unconditioned generation plus post-hoc filtering, an
oracle finite-catalogue assembler, a learned inventory selector, two learned component-factorized
generators, and targeted interventions on the Transformer. Appendix~\ref{app:ablation-results}
reports the three-seed Transformer mechanism and component-factorization study. Each
intervention used the same folds, generation budget and evaluation code as the full Transformer.

\subsubsection{Metrics, molecular identity and uncertainty}

Let $Y_{s,a,b}^{\rm L1}\in\{0,1\}$ indicate verified exact L1 for training seed $s$, arm $a$ and
attempt $b$, and let $B=\ForgeProseAttemptsPerProgram$. We report the exact count
$C_{s,a}^{\rm L1}=\sum_{b=1}^B Y_{s,a,b}^{\rm L1}$ and
\begin{equation}
 \widehat p_{s,a}^{\rm L1}=C_{s,a}^{\rm L1}/B.
\end{equation}
Validity and connected yield replace $Y^{\rm L1}$ with their attempt-level indicators. Verified
decomposition coverage is the number of valid graphs with
$\widehat{\mathcal E}_c(G)\neq\varnothing$ divided by valid graphs; the implementation does not
report candidate-search coverage before replay. Exact ambiguity is the fraction of valid graphs with
$|\widehat{\mathcal E}_c(G)|>1$. Replay of members of $\widehat{\mathcal E}_c(G)$ is an invariant,
whereas candidate replay precision is defined in Appendix~\ref{app:verification}. The evaluated
artifacts do not retain the pre-replay candidate sets, so candidate-search coverage and candidate
replay precision were not estimated.
Explicitly, with attempt outcomes $G_{s,a,b}$ and the convention that a failed decode makes every
graph predicate zero,
\begin{align}
 \widehat p^{\rm valid}_{s,a}
 &=B^{-1}\sum_b\mathbf 1\{\mathrm{Valid}(G_{s,a,b})=1\},\\
 \widehat p^{\rm conn}_{s,a}
 &=B^{-1}\sum_b\mathbf 1\{\mathrm{Valid}(G_{s,a,b})=1,
                              \mathrm{Connected}(G_{s,a,b})=1\},\\
 \mathrm{CandidateCoverage}
 &=\frac{\sum_b\mathbf 1\{\mathrm{Valid}(G_b)=1,
                              \mathcal D_c(G_b)\neq\varnothing\}}
         {\sum_b\mathbf 1\{\mathrm{Valid}(G_b)=1\}},\\
 \mathrm{ExactAmbiguity}
 &=\frac{\sum_b\mathbf 1\{\mathrm{Valid}(G_b)=1,
                              |\widehat{\mathcal E}_c(G_b)|>1\}}
         {\sum_b\mathbf 1\{\mathrm{Valid}(G_b)=1\}}.
\end{align}
Each ratio is N/E when its denominator is zero.

Distinct-product yield counts canonical constitutions once. The reported verifier-recovered
component-novel yield also counts distinct products and requires an unambiguous recovered input
with an out-of-training-catalogue precursor. It does not measure strong catalogue escape
(Definition~\ref{def:catalogue-escape}). Held-component recovery counts qualifying attempts,
including duplicate products, whose unambiguous input contains a designated held constitution.

For one chemistry $c$, suppress seed and arm indices. Let $\mathcal A_{\rm L1}$ contain the attempt indices
with valid graphs and verified exact L1, and let
$\mathcal A_{\rm one}=\{b\in\mathcal A_{\rm L1}:|\widehat{\mathcal E}_c(G_b)|=1\}$.
For $b\in\mathcal A_{\rm one}$, write the sole recovered input as $\mathbf C_b$, with components $C_{b,j}$
and roles $r_{b,j}$. Let $\mathcal C_{{\rm train},r}$ and $\mathcal H_{{\rm held},r}$ be the
training and designated held constitution sets for role $r$, respectively. Define
\begin{align}
 \mathcal A_{\rm novel}&=\{b\in\mathcal A_{\rm one}:\exists j,\ C_{b,j}\notin\mathcal C_{{\rm train},r_{b,j}}\},\\
 \mathcal A_{\rm held}&=\{b\in\mathcal A_{\rm one}:\exists j,\ C_{b,j}\in\mathcal H_{{\rm held},r_{b,j}}\}.
\end{align}
With $\operatorname{can}$ denoting constitutional canonicalization, the yields for $B>0$ are
\begin{align}
 \mathrm{DistinctL1Yield}&=\frac{|\{\operatorname{can}(G_b):b\in\mathcal A_{\rm L1}\}|}{B},\\
 \mathrm{ComponentNovelYield}&=\frac{|\{\operatorname{can}(G_b):b\in\mathcal A_{\rm novel}\}|}{B},\\
 \mathrm{HeldComponentYield}&=\frac{|\mathcal A_{\rm held}|}{B}.
\end{align}
Reported values per 1,000 attempts multiply these ratios by 1,000. Ambiguous decompositions do not
contribute to the component-novel or held-component numerators.

For a fixed role, let $n_i$ be the occurrence count of component identity $i$ among unambiguous
exact-L1 inputs. When $n_{\rm tot}=\sum_i n_i>0$, set $w_i=n_i/n_{\rm tot}$ and compute
\begin{equation}
 N_{\rm eff}=\exp\!\left(-\sum_{i:n_i>0}w_i\log w_i\right).
\end{equation}
This is the exponential of Shannon entropy with natural logarithms. Duplicate products and repeated
component occurrences contribute to the counts; $N_{\rm eff}$ is N/E when $n_{\rm tot}=0$.
Molecular diversity is the mean pairwise
ECFP4 Tanimoto distance among distinct valid canonical product constitutions, so molecular duplicates
are removed for that metric. For distinct products $g_1,\ldots,g_m$,
\begin{equation}
 \mathrm{Diversity}=
 \frac{2}{m(m-1)}\sum_{1\leq i<j\leq m}
 \left[1-\mathrm{Tanimoto}\bigl(\mathrm{ECFP4}(g_i),\mathrm{ECFP4}(g_j)\bigr)\right],
\end{equation}
with N/E for $m<2$. Every other undefined conditional denominator is likewise reported N/E.

Identity is constitutional and stereo-free. Graph and product identity use connected canonical RDKit
SMILES with atom maps and stereochemical tags removed; formal charges and bond orders are retained;
salts or disconnected products are invalid; and no tautomer normalization is applied. Precursor
identity uses the same canonicalization. Precursor-tuple equality is role- and program-step-aware,
with repeated inputs compared in their registered step order. Training atom mappings and recovered assembly traces serve different roles. The former establish
semantic alignment; the latter record precursor assignments and reaction steps. Neither is part
of constitutional identity.

For a paired contrast, $D_s=\widehat p_{s,\rm FORGE}-\widehat p_{s,\rm control}$. We report
$(D_1,D_2,D_3)$, $\overline D$, and $[\min_sD_s,\max_sD_s]$. These summarize the three observed
training seeds and are not 95\% confidence intervals. Conditional on one fixed trained model,
$C_{s,a}^{\rm L1}$ has the finite-attempt binomial variation in
Proposition~\ref{prop:posthoc-budget}; that uncertainty does not substitute for independent
retraining. Unless explicitly named primary, architecture, coordinate, guidance, potency and
catalogue contrasts are exploratory and no multiplicity-adjusted inferential claim is made.

\FloatBarrier
\subsection{Assembly quality and program conditioning}
\label{app:assembly-results}

\subsubsection{Shared generation across three chemistries}

The full shared-model comparison used \ForgeProseAttemptsPerProgram{} held-out generation attempts
per program in each of three independent training seeds (Table~\ref{tab:production-comparison}).
Table~\ref{tab:iclr-shared-programs} summarizes the exact-L1 yields.
Every verifier-returned trace passed forward replay, as required by the verifier invariant. Mean exact-
decomposition ambiguity was \ForgeTableOneBLAmbiguityPercent\% for repeated aza-Michael addition
and \ForgeTableOneLXAmbiguityPercent\% for
repeated reductive amination, and the latter program showed substantial between-seed variation.
Table~\ref{tab:production-comparison} reports molecular validity, connectedness, decomposition
coverage, forward-replay precision, abstention, ambiguity, novelty and diversity separately for
each program.

The final conditioned model recorded zero fixed-coordinate violations and zero support overflows across
\ForgeProseConditionedTotalAttempts{} held-out attempts (three seeds, three programs and
\ForgeProseAttemptsPerProgram{} attempts per program).
Exact-assembly yields differed across reaction families despite preservation of all fixed
coordinates and graph-support bounds.

% The compact shared-program table is in the main paper as Table~\ref{tab:iclr-shared-programs}.

\begin{table*}[htbp]
\caption{\textbf{Matched shared-model production comparison across reaction programs.}
Point estimates report the mean across three independent seeds; the exact-L1 column additionally
reports the sample standard deviation. Paired contrasts are summarized by their three seed values,
mean and observed range. Conditioned-model counts by seed appear in
Table~\ref{tab:production-seed-counts}.
Validity and connectedness use all generation attempts as the denominator.
Decomposition coverage, replay precision, abstention and ambiguity retain their prespecified
conditional denominators. Novelty is conditioned on valid products or unique returned
decompositions, as appropriate. Fraction pairs and exact-L1 yield are percentages; diversity is
mean pairwise ECFP4 distance and $N_{\mathrm{eff}}$ is the effective generated-component count.
Results are never pooled across reaction programs.}
\label{tab:production-comparison}
\centering
\scriptsize
\setlength{\tabcolsep}{3.5pt}
\resizebox{\textwidth}{!}{%
\begin{tabular}{@{}lccccccccc@{}}
\hline
& \multicolumn{3}{c}{Ugi} & \multicolumn{3}{c}{Aza-Michael}
& \multicolumn{3}{c}{Reductive amination} \\
\cmidrule(lr){2-4}\cmidrule(lr){5-7}\cmidrule(lr){8-10}
& FORGE & Shared null & Cyclic & FORGE & Shared null & Cyclic
& FORGE & Shared null & Cyclic \\
\hline
\input{generated/production_comparison_transposed_rows.tex}%
\end{tabular}
}
\end{table*}



\begin{table}[htbp]
\caption{\textbf{Exact verified-L1 counts for the final conditioned model.}
Each row reports the numerator, fixed attempt denominator and resulting percentage for one
independently trained seed. These attempt-level counts quantify conditional Monte Carlo variation;
the three training seeds are the independent units for retraining variation. Seed
labels 0, 1 and 2 correspond to RNG seeds 20260825, 20260826 and 20260827, respectively.}
\label{tab:production-seed-counts}
\centering
\small
\setlength{\tabcolsep}{7pt}
\begin{tabular}{lrrrr}
\hline
Program & Seed & Verified L1 count & Attempts & Yield (\%) \\
\hline
\input{generated/production_seed_exact_counts_rows.tex}%
\end{tabular}
\end{table}

\subsubsection{Conditioned, null and cyclic controls}

Under the matched protocol in Table~\ref{tab:production-comparison}, the paired
conditioned-minus-null differences were
\ForgeTableOneNullUgiDifferencePP{} percentage points for Ugi (three-seed range
\ForgeTableOneNullUgiRangeLowPP{} to \ForgeTableOneNullUgiRangeHighPP{}),
\ForgeTableOneNullBLDifferencePP{} for aza-Michael (\ForgeTableOneNullBLRangeLowPP{} to
\ForgeTableOneNullBLRangeHighPP{}), and \ForgeTableOneNullLXDifferencePP{} for reductive amination
(\ForgeTableOneNullLXRangeLowPP{} to \ForgeTableOneNullLXRangeHighPP{}).

The cyclic control retained precursor roles, reaction-core coordinates and program depth but assigned
the wrong reaction identity.
Relative to correct conditioning, this intervention changed verified exact-L1 yield by
\ForgeTableOneCyclicUgiDifferencePP{} percentage points for Ugi (three-seed range
\ForgeTableOneCyclicUgiRangeLowPP{} to \ForgeTableOneCyclicUgiRangeHighPP{}),
\ForgeTableOneCyclicBLDifferencePP{} for aza-Michael (\ForgeTableOneCyclicBLRangeLowPP{} to
\ForgeTableOneCyclicBLRangeHighPP{}), and \ForgeTableOneCyclicLXDifferencePP{} for reductive
amination (\ForgeTableOneCyclicLXRangeLowPP{} to \ForgeTableOneCyclicLXRangeHighPP{}). Every
returned exact trace passed replay by verifier construction.

We independently retrained the null and cyclic arms under the same study protocol to measure how
program information affected learned generation. Correct conditioning improved yield substantially
relative to the null arm. Its smaller advantage over cyclic labels did not establish a strong
independent contribution from the reaction-identity token.

\FloatBarrier
\subsection{Decoder and architecture ablations}
\label{app:ablation-results}

\subsubsection{Terminal decoding and source marginals}

We evaluated terminal decoding and source marginals in separate fixed-weight comparisons
(Table~\ref{tab:decoder-source-ablation}).

\begin{table}[htbp]
\caption{\textbf{Decoder and source-marginal ablations of the final model.} Ugi held-out attempts
for the 9,143-step shared production model. The two
ablations are independent: the decoder rows share one set of trained weights and differ only in
terminal decoding, while the source rows share one decoder and differ only in the marginal the
corruption starts from. Decoding under the reaction-core saturation policy holds the aldehyde-derived
$\alpha$-carbon at one hydrogen and the isocyanide-derived amide nitrogen at one hydrogen, which the
qualified transform fixes; it changes the generated graph under an unchanged verifier. Values
are percentages from \ForgeProseAttemptsPerProgram{} seed-0 held-out attempts per row; they are descriptive and do not measure
across-training-seed variation. Exact-L1 coverage is conditioned on valid products.}
\label{tab:decoder-source-ablation}
\centering
\scriptsize
\setlength{\tabcolsep}{5pt}
\begin{tabular}{@{}llccc@{}}
\toprule
Source marginal & Terminal decoder & Valid & Exact-L1 coverage & Exact-L1 \\
\midrule
\input{generated/decoder_source_ablation_rows.tex}%
\end{tabular}
\end{table}

\subsubsection{Joint denoising, role isolation and Transformer modules}

The matched-study full Transformer produced \ForgeProseMechanismFullUgi{} Ugi and
\ForgeProseMechanismFullBL{} repeated aza-Michael exact-L1 products per 1,000 attempts, compared with
\ForgeProseMechanismFactMatchedUgi{} and \ForgeProseMechanismFactMatchedBL{} for the
parameter-matched role-isolated FACT model
(Table~\ref{tab:architecture-ablations}). The full Transformer and FACT-matched model each had
\ForgeProseMechanismSharedParameterMillions{} million trainable parameters; FACT-generous had
\ForgeProseMechanismGenerousParameterMillions{} million. The higher-capacity FACT control improved
to \ForgeProseMechanismFactGenerousUgi{} for Ugi and \ForgeProseMechanismFactGenerousBL{} for
repeated aza-Michael addition. Its three-seed means remained
below the full Transformer for both programs, although the aza-Michael values were essentially tied
in one seed. The direction reversed for repeated reductive amination: the full Transformer produced
\ForgeProseMechanismFullLX{} exact-L1 products per 1,000 attempts, whereas FACT-matched and
FACT-generous produced \ForgeProseMechanismFactMatchedLX{} and
\ForgeProseMechanismFactGenerousLX{}. Thus, joint whole-lipid denoising improved exact assembly
for Ugi and repeated aza-Michael addition but was not uniformly superior across reaction programs.
Because the prespecified conditional-dependence diagnostic was negative, this architecture contrast
does not establish detectable cross-role dependence in the present data.

\begin{table*}[htbp]
\caption{\textbf{Transformer architecture and mechanism ablations.}
The five Transformer intervention rows change only the named mechanism from the matched-study full
model; the two
FACT rows generate precursor roles independently at matched or increased capacity. All arms
use identical data folds, corruption, context schedule, checkpoint rule, core-saturation terminal
decoder and sampling budget. The reference full Transformer is the matched mechanism-study control,
not the independently trained production model reported in Table~\ref{tab:production-comparison}
and Appendix~\ref{app:novelty-results}. Entries
report the mean $\pm$ sample standard deviation across three independent training seeds.
Held-out loss is the mean denoising loss at $t=0.5$ across the three component-disjoint program partitions.
Cross-role fidelity is the Frobenius distance to the held-out residualized descriptor matrix, for
which lower values indicate closer agreement. Exact-L1 yields use all
\ForgeProseAttemptsPerProgram{} attempts per program
and seed as the denominator.}
\label{tab:architecture-ablations}
\centering
\scriptsize
\setlength{\tabcolsep}{2.0pt}
\resizebox{\textwidth}{!}{%
\begin{tabular}{@{}>{\raggedright\arraybackslash}p{2.4cm}ccccccc@{}}
\hline
Arm &
Held-out loss &
Ugi L1/1k &
BL L1/1k &
LX L1/1k &
Held-comp. L1/1k &
Cross-role fidelity &
Diversity \\
\hline
\input{generated/architecture_ablation_rows.tex}%
\end{tabular}
}
\end{table*}

Removing layerwise program cross-attention reduced mean exact-L1 yield by
\ForgeProseNoCrossAttentionUgiReduction{} products per 1,000 attempts for Ugi and
\ForgeProseNoCrossAttentionLXReduction{} for repeated reductive amination, while increasing repeated
aza-Michael yield by \ForgeProseNoCrossAttentionBLIncrease{}. Removing the role or core objective and
removing routed adapters also produced
program-dependent changes. The arm without gradient-conflict control slightly exceeded the full
Transformer mean in all three programs. These interventions did not support a uniform benefit
from every module. We reported effects separately by chemistry and did not combine them into an
architecture score.

\FloatBarrier
\subsection{Component novelty, generalization and structural realism}
\label{app:novelty-results}

\subsubsection{Common Ugi benchmark}

\label{app:reporting}

We compared FORGE with external generators, matched internal controls and the finite-catalogue
oracle under the common Ugi benchmark (Table~\ref{tab:ugi-common-benchmark}).
Every yield is normalized by the number of generation attempts, including attempts that fail before
producing a valid molecule.
All included methods were run under the common split, attempt budget, seed set and verifier.

% float moved to the main body: tab:ugi-common-benchmark

The learned finite-inventory selector achieved \ForgeProseCommonSelectorExactPerThousand{} verified
exact-L1 products per 1,000
attempts, but its finite action space imposed zero verifier-recovered component-novel yield. FORGE
produced \ForgeProseCommonForgeExactPerThousand{} verified exact-L1 products and
\ForgeProseCommonForgeDistinctPerThousand{} distinct verified exact-L1 products per 1,000 attempts.
Among whole-molecule external generators, unconditional DeFoG produced
\ForgeProseCommonDefogExactPerThousand{} verified exact-L1 products per 1,000 attempts and recovered
held components at \ForgeProseCommonDefogHeldPerThousand{} per 1,000 attempts, compared with
\ForgeProseCommonForgeHeldPerThousand{} for FORGE. GenMol/SAFE produced
\ForgeProseCommonGenmolValidPerThousand{} valid connected outputs but no exact Ugi-L1 product, while
RGFN emitted no admissible
molecular output under the strict native contract. These results established different performance orderings for
reaction consistency, verifier-recovered component novelty and recovery of reserved components. They do not establish product-level strong catalogue escape.

Seed-level results and decomposition outcomes are reported in Tables~\ref{tab:baseline-seed-results}
and~\ref{tab:baseline-decomposition}. All values derive from hash-pinned three-seed results.
Compute records are
retained in the generated evidence receipt, but missing training-time fields prevent a
cross-method compute-matching claim.

\begin{table}[htbp]
\caption{\textbf{Seed-level common Ugi benchmark results.}
Each method has a hash-verified three-seed result vector. Table~\ref{tab:ugi-common-benchmark}
reports the corresponding mean $\pm$ sample standard deviation. Training seeds are the independent
units; generated molecules are samples from each fixed model.}
\label{tab:baseline-seed-results}
\centering
\scriptsize
\setlength{\tabcolsep}{3.5pt}
\resizebox{\textwidth}{!}{%
\begin{tabular}{llcccccc}
\hline
Method & Seed(s) & Valid/1k & Verified L1/1k & Distinct L1/1k & Comp.-novel/1k & Held-comp. L1/1k & Diversity \\
\hline
\input{generated/common_ugi_seed_rows.tex}%
\end{tabular}
}
\end{table}

\begin{table}[htbp]
\caption{\textbf{Reaction decomposition outcomes under the common Ugi verifier.}
Verified decomposition coverage is conditioned on molecular validity. Replay of every tuple in the
returned exact set is a verifier invariant, not an estimate of reverse-search precision. Abstention
and exact ambiguity retain their prespecified denominators. N/E denotes a conditional metric with
no eligible valid product or returned decomposition.}
\label{tab:baseline-decomposition}
\centering
\scriptsize
\setlength{\tabcolsep}{4.5pt}
\resizebox{\textwidth}{!}{%
\begin{tabular}{lccccc}
\hline
Method & Valid products & Verified coverage & Replay invariant & Abstention & Exact ambiguity \\
\hline
\input{generated/common_ugi_decomposition_rows.tex}%
\end{tabular}
}
\end{table}

\subsubsection{Catalogue comparison and component diversity}

At the matched sampling budget, FORGE produced
\ForgeProseCatalogueForgeUgiComponentNovelPerThousand{} verifier-recovered component-novel Ugi
products per 1,000 attempts, averaged across the three independent seeds. The corresponding
values were \ForgeProseCatalogueForgeBLComponentNovelPerThousand{} for aza-Michael addition and
\ForgeProseCatalogueForgeLXComponentNovelPerThousand{} for reductive amination. The finite-component
assembler produced zero by construction because every sampled product uses only components from
the training catalogue (Table~\ref{tab:catalogue-comparison}).
The catalogue achieved exact-L1 yields of
\ForgeProseCatalogueCatalogueUgiExactPerThousand{},
\ForgeProseCatalogueCatalogueBLExactPerThousand{} and
\ForgeProseCatalogueCatalogueLXExactPerThousand{} products per 1,000 attempts, compared with
\ForgeProseCatalogueForgeUgiExactPerThousand{},
\ForgeProseCatalogueForgeBLExactPerThousand{} and
\ForgeProseCatalogueForgeLXExactPerThousand{} for FORGE. The corresponding raw validities were
\ForgeProseCatalogueCatalogueUgiValidPerThousand{},
\ForgeProseCatalogueCatalogueBLValidPerThousand{} and
\ForgeProseCatalogueCatalogueLXValidPerThousand{} for the catalogue and
\ForgeProseCatalogueForgeUgiValidPerThousand{},
\ForgeProseCatalogueForgeBLValidPerThousand{} and
\ForgeProseCatalogueForgeLXValidPerThousand{} for FORGE.
We report these quantities separately because reaction validity and verifier-relative component
novelty measure different capabilities. Direct membership in the complete catalogue-reachable
product set was not computed, so these results are not reported as strong catalogue escape.

In the distinct-product Ugi role analysis, \ForgeOffRegistryCount{} of
\ForgeOffRegistryDenominator{} products (\ForgeOffRegistryPercent\%) contained at least one
verifier-returned component constitution absent from the training registry. Across all Ugi products
included in the novelty analysis, \ForgeNovelAgainstFullCount{} of
\ForgeNovelAgainstFullDenominator{} distinct products (\ForgeNovelAgainstFullPercent\%) were absent
from the full \ForgeProseFullReferenceSize{}-product constitutional reference corpus. FORGE effective
component counts were \ForgeProseCatalogueForgeUgiEffectiveComponentCount{},
\ForgeProseCatalogueForgeBLEffectiveComponentCount{} and
\ForgeProseCatalogueForgeLXEffectiveComponentCount{}, compared with
\ForgeProseCatalogueCatalogueUgiEffectiveComponentCount{},
\ForgeProseCatalogueCatalogueBLEffectiveComponentCount{} and
\ForgeProseCatalogueCatalogueLXEffectiveComponentCount{} for the finite catalogue. The Ugi catalogue
contained \ForgeProseCatalogueUgiAmineCount{} amines,
\ForgeProseCatalogueUgiAldehydeCount{} aldehydes and
\ForgeProseCatalogueUgiIsocyanideCount{} isocyanides and covered
$\ForgeProseCatalogueTrainCoveragePercent\%$ of training products but
$\ForgeProseCatalogueHeldCoveragePercent\%$ of the calibration and component-disjoint held-out
products. Its Cartesian tuple ceilings were \ForgeProseCatalogueUgiTupleCeiling{} for Ugi,
\ForgeProseCatalogueBLTupleCeiling{} for aza-Michael and
\ForgeProseCatalogueLXTupleCeiling{} for reductive amination; these count possible
component-program inputs, whose forward products may overlap. The catalogue comparison therefore
quantified both assembly consistency and the expansion of recovered precursor identities beyond
that inventory.

\begin{table}[htbp]
\caption{\textbf{Whole-graph generation versus a finite-component assembler.}
Verified exact-L1, distinct exact-L1 and verifier-recovered component-novel yields are reported per
1,000 generation attempts. A component-novel product has exactly one verifier-returned exact-L1
decomposition and contains at least one component constitution absent from the training catalogue.
This verifier-relative metric is not strong catalogue escape because complete catalogue-product
reachability was not computed. The catalogue has zero input-component novelty by construction, but may
achieve higher raw validity or exact-L1 yield. FORGE-minus-catalogue contrasts use the three seeds
as descriptive paired units; no three-seed confidence interval is inferred.}
\label{tab:catalogue-comparison}
\centering
\scriptsize
\setlength{\tabcolsep}{4pt}
\resizebox{\textwidth}{!}{%
\begin{tabular}{@{}lcccccc@{}}
\hline
& \multicolumn{2}{c}{Ugi} & \multicolumn{2}{c}{Aza-Michael}
& \multicolumn{2}{c}{Reductive amination} \\
\cmidrule(lr){2-3}\cmidrule(lr){4-5}\cmidrule(lr){6-7}
& FORGE & Finite cat. & FORGE & Finite cat. & FORGE & Finite cat. \\
\hline
\input{generated/catalogue_comparison_transposed_rows.tex}%
\end{tabular}
}
\end{table}

\subsubsection{Held-component recovery}

The component-disjoint Ugi fold was never used for training, checkpoint selection or inventory
construction. Under the common native-generation assessment, the full Transformer recovered only
\ForgeProseCommonHeldTotal{} exact-L1 products containing a designated held component in
\ForgeProseCommonHeldAttempts{} attempts, or \ForgeProseCommonForgeHeldPerThousand{} products per 1,000
attempts averaged across three seeds. Recovery of specifically reserved components was therefore weak despite the broader component
novelty observed above. Novelty accepts any recovered precursor absent from training; held-component
recovery requires a match to a designated withheld identity. The results do not establish strong
held-component recovery, role-specific held-component coverage or zero-shot generalization to
held reaction families.

\subsubsection{Structural realism}

\begin{table*}[htbp]
\caption{\textbf{Method-blind structural-realism diagnostic.}
Every entry is mean $\pm$ sample standard deviation across three training seeds and retains all
\ForgeProseAttemptsPerProgram{} requested attempts in the denominator. FP and descriptor P/C report manifold precision per
1,000 attempts and coverage of the source-study-held-out R0 reference in percent. C2ST is grouped
classifier two-sample area under the receiver operating characteristic curve; 0.5 denotes
chance discrimination. $N_{\mathrm{eff}}$ is the Shannon effective molecule count. N/E denotes a metric that cannot be
estimated. This post-hoc structural diagnostic is not a property, activity or synthesis-success
score.}
\label{tab:lipid-realism}
\centering
\scriptsize
\setlength{\tabcolsep}{1.6pt}
\resizebox{\textwidth}{!}{%
\begin{tabular}{@{}>{\raggedright\arraybackslash}p{2.4cm}cccccc@{}}
\hline
Method & Connected/1k & FP P/C & Desc. P/C & C2ST & $N_{\mathrm{eff}}$ & Diversity \\
\hline
\input{generated/lipid_realism_rows.tex}%
\end{tabular}
}
\end{table*}

No evaluated method reproduced the full held-out observed-lipid distribution: where estimable,
grouped C2ST AUC remained between \ForgeProseRealismCtwoSTMin{} and
\ForgeProseRealismCtwoSTMax{}, and fingerprint-manifold precision was zero or
near zero (Table~\ref{tab:lipid-realism}). GenMol/SAFE's high pairwise diversity was paired with an
effective count of only \ForgeProseRealismGenmolEffective{} molecules, confirming severe frequency
collapse. FORGE retained \ForgeProseRealismForgeEffective{} effective molecules, but only
\ForgeProseRealismForgeDescriptorPrecision{} attempts per 1,000 fell within the descriptor
manifold. Diversity, effective molecule count and reference-distribution coverage therefore
measured different aspects of the generated populations.

% float moved to the main body: fig:forge-generated-samples

\FloatBarrier
\subsection{Potency-guidance diagnostic}
\label{app:guidance-results}

We evaluated biological guidance only in the bounded HeLa messenger-transfection-potential (mTP)
diagnostic. A group-disjoint preterminal classifier detected weak morphology-level signal, with ROC
AUC \ForgePotencyAUC{} (95\% interval \ForgePotencyAUCCILow{} to \ForgePotencyAUCCIHigh) and a
\ForgePotencyTopQuartileGain-fold observed top-quartile gain. We then compared a frozen
applicability-enriched morphology proposal with a nested potency-enriched proposal using
\ForgePotencyAttemptsPerArm{} generator calls per arm. Table~\ref{tab:hela-property-guidance}
reports the matched terminal outcome.

\begin{table}[htbp]
\caption{\textbf{Predicted HeLa messenger-transfection potential under matched generation budgets.}
The support-enriched arm used the selected applicability proposal. The nested arm reallocated
morphology mass using the frozen weak potency signal. Conservative-high counts are computational
predictions, not transfection measurements. The comparison uses the same generator-call and oracle
budgets.}
\label{tab:hela-property-guidance}
\centering
\scriptsize
\setlength{\tabcolsep}{5.0pt}
\resizebox{\textwidth}{!}{%
\begin{tabular}{lccccc}
\hline
Arm &
Attempts &
Verified L1 &
Eligible before budget &
Oracle budget &
Unique conservative-high \\
\hline
Support-enriched & \ForgePotencyAttemptsPerArm{} & \ForgePotencySupportExactLone{} &
\ForgePotencySupportEligible{} & \ForgePotencyOracleBudget{} & \ForgePotencySupportHigh{} \\
Nested potency proposal & \ForgePotencyAttemptsPerArm{} & \ForgePotencyNestedExactLone{} &
\ForgePotencyNestedEligible{} & \ForgePotencyOracleBudget{} & \ForgePotencyNestedHigh{} \\
\hline
\end{tabular}
}
\end{table}

The nested-minus-support difference was \ForgePotencyDifferencePerThousand{} conservative-high
products per 1,000 generator calls (95\% interval \ForgePotencyDifferenceCILowPerThousand{} to
\ForgePotencyDifferenceCIHighPerThousand). The comparison failed the prespecified advancement criterion. These computations evaluate
predictor behavior within a declared applicability domain; they do not establish improved biological
activity, delivery or experimental potency.


\clearpage

\section{Theoretical analysis and proofs}
\label{app:theory}

The results establish separate properties of the construction. Bounded completeness concerns the
state representation. Catalogue reachability concerns exact precursor inputs. The population
argument connects a reference denoising loss to an ideal continuous-time law. Budget and route
results concern computational procedures. None establishes experimental synthesis success or
positive learned probability for every representable graph.

\subsection{Graph states and program compatibility}
\label{app:state-support}

The graph space $\mathfrak G_{N,K_{\rm cl}}$ is defined in
Section~\ref{sec:mathematical-objects}, with finite atom and bond alphabets and integers $N\geq1$,
$K_{\rm cl}\geq0$. Molecular validity is an additional predicate on that space.
The raw space $\overline{\mathfrak X}_{N,K_{\rm cl}}$ contains padded categorical states.

\begin{definition}[Well-formed sparse state]
\label{def:sparse-state}
A state belongs to $\mathfrak X_{N,K_{\rm cl}}$ when, for some $1\leq n\leq N$, its active atoms
are exactly positions $1,\ldots,n$ with labels in $\Sigma_{\rm atom}$; position 1 is the sole root
with null parent and parent-bond states; and every $2\leq i\leq n$ has an earlier parent
$1\leq\pi_i<i$ and a bond label in $\Sigma_{\rm bond}$. Each active closure has endpoints
$1\leq e^L<e^R\leq n$, a bond label in $\Sigma_{\rm bond}$, and an edge distinct from every tree
edge and other closure. Active closures precede unused slots and are sorted by endpoint pair.
All unused atom, pointer and bond coordinates carry their designated null values.
\end{definition}
The decoder reconstructs the labeled edges of a well-formed state and returns $\bot_{\rm dec}$
otherwise. It is therefore a total map into $\mathfrak G_{N,K_{\rm cl}}\cup\{\bot_{\rm dec}\}$.

\begin{lemma}[Tree residual count]
\label{lem:tree-closure}
For a finite nonempty connected graph $G=(V,E)$ and any spanning tree $T$,
$|E\setminus E(T)|=|E|-|V|+1=\beta(G)$.
\end{lemma}
\begin{proof}
The tree contains $|V|-1$ edges, all in $E$. Subtracting this count from $|E|$ gives the result,
including the single-vertex case.
\end{proof}

\begin{theorem}[Bounded completeness of the sparse representation]
\label{thm:bounded-representation}
For the spaces in Definition~\ref{def:sparse-state}, the restriction
$\operatorname{Dec}:\mathfrak X_{N,K_{\rm cl}}\to\mathfrak G_{N,K_{\rm cl}}$ is well defined
and surjective. Thus every well-formed state decodes to a graph in the declared support, and every
graph in that support has an encoding that decodes to its constitutional isomorphism class.
\end{theorem}
\begin{proof}
We first reconstruct a graph from a state, then construct a state for an arbitrary graph.
In a well-formed state with $n$ active atoms, every parent pointer strictly decreases its index.
Following parents must reach position 1, so the $n-1$ parent edges form a connected tree.
Adding $k\leq K_{\rm cl}$ distinct, non-loop closure edges preserves connectedness and gives cycle
rank $(n-1+k)-n+1=k$. The label and padding conditions ensure the decoded graph belongs to
$\mathfrak G_{N,K_{\rm cl}}$.

Conversely, choose a representative of any $G\in\mathfrak G_{N,K_{\rm cl}}$ and a rooted spanning
tree. Enumerate vertices in a traversal that places every parent before its children. Store their
labels, parent indices and tree-bond labels. Lemma~\ref{lem:tree-closure} leaves exactly
$\beta(G)\leq K_{\rm cl}$ residual edges; store each with increasing endpoint indices and sort
these triples. Pad all unused coordinates with null states. This state satisfies
Definition~\ref{def:sparse-state}, and the vertex enumeration is a label-preserving isomorphism
between its decoded graph and $G$.
\end{proof}

This result concerns the unrestricted representation. For program-specific alphabets
$\mathcal X_i(P)$ and fixed values $x_{P,i}$, define
\begin{equation}
 \Omega_P=\{x\in\overline{\mathfrak X}_{N,K_{\rm cl}}:
 x_i\in\mathcal X_i(P)\ \forall i,\quad x_i=x_{P,i}\ \forall i\in\mathcal F(P)\}.
 \label{eq:program-state-set}
\end{equation}
Its nonempty well-formed intersection defines
$\mathcal P_{\rm adm}=\{P:\mathfrak X_{N,K_{\rm cl}}\cap\Omega_P\neq\varnothing\}$.
Nonemptiness alone establishes neither chemical validity nor satisfaction of every morphology
constraint. Corpus admission additionally checks the annotated encoding, program alignment and
exact decode round trip. Representational completeness makes no claim that every graph is
encodable under every $P$.

\subsection{Exact assembly and catalogue reachability}
\label{app:catalogue-reachability}

Definition~\ref{def:exact-assembly} distinguishes the complete input preimage
$\mathcal E_c^\star(G)$ from recovered input--trace pairs $\widehat{\mathcal E}_c(G)$.
Fix a registered chemistry $c$ and role-specific catalogues $\mathcal C=(\mathcal C_r)_r$. Let
$\mathfrak C_c(\mathcal C)\subseteq\mathfrak C_c$ contain the admissible precursor inputs whose
every role-assigned component belongs to its catalogue. The reachable product set is
\begin{equation}
 \mathcal G_{\rm cat}(c,\mathcal C)
 =\bigcup_{\mathbf C\in\mathfrak C_c(\mathcal C)}\mathcal F_c(\mathbf C).
 \label{eq:catalogue-reachable}
\end{equation}
This definition includes all registered outcomes of the admissible inputs, not only products found
by a bounded search.

\begin{proposition}[Finite-component reachability]
\label{prop:catalogue-bound}
If $\mathfrak C_c(\mathcal C)$ and each $\mathcal F_c(\mathbf C)$ are finite, then
\begin{equation}
 |\mathcal G_{\rm cat}(c,\mathcal C)|
 \leq\sum_{\mathbf C\in\mathfrak C_c(\mathcal C)}|\mathcal F_c(\mathbf C)|
 \leq m_{\max,c}(\mathcal C)\,|\mathfrak C_c(\mathcal C)|,
\end{equation}
where $m_{\max,c}(\mathcal C)$ is the largest outcome count, or zero for an empty input set.
For every $G\in\mathfrak G_{N,K_{\rm cl}}$,
\begin{equation}
 G\in\mathcal G_{\rm cat}(c,\mathcal C)
 \quad\Longleftrightarrow\quad
 \mathcal E_c^\star(G)\cap\mathfrak C_c(\mathcal C)\neq\varnothing.
\end{equation}
\end{proposition}
\begin{proof}
The cardinality of a finite union is at most the sum of its member-set cardinalities; bounding each
summand by its maximum proves both inequalities. If there are no inputs, all three quantities are
zero. Membership in the union means that some catalogue input has $G$ as a forward product, which
is exactly the stated intersection condition.
\end{proof}

\begin{definition}[Strong catalogue escape]
\label{def:catalogue-escape}
An exactly assemblable graph strongly escapes $\mathcal C$ under $c$ when every exact input uses
at least one out-of-catalogue component. Its indicator is
\begin{equation}
 \operatorname{Escape}^\star(G;c,\mathcal C)
 =\mathbf1\!\left[\mathcal E_c^\star(G)\neq\varnothing,\quad
 \mathcal E_c^\star(G)\cap\mathfrak C_c(\mathcal C)=\varnothing\right].
\end{equation}
\end{definition}
A catalogue assembler whose outputs lie in $\mathcal G_{\rm cat}(c,\mathcal C)$ therefore has
zero strong-escape yield. For a verified exact-L1 product, strong escape is equivalent to
$G\notin\mathcal G_{\rm cat}(c,\mathcal C)$. This can be checked against a complete canonical
product set or by exhaustive exclusion of all catalogue-only decompositions. A bounded reverse
search cannot establish that exclusion from one novel witness.

The reported experiments did not construct a complete common-reference reachable product set.
Their component-novelty metric concerns the selected verified input. Whole-graph generation removes
catalogue membership from the generative representation; that architectural property alone proves
neither strong escape nor positive probability of any particular out-of-catalogue graph.

\subsection{Conditional flow and ideal population recovery}
\label{app:flow-proof}

We prove Theorem~\ref{thm:ideal-terminal-law} for the reference objective in
Equation~\ref{eq:primary-loss}. Let $X_\star=\operatorname{Enc}_P(G,\alpha)$ denote a clean training
endpoint, and distinguish it from the time-indexed process $X_t$. The training-path measure is
\begin{equation}
 \nu(dG,d\alpha,dP,dt,dx)
 =\mu_{\rm bal}(dG,d\alpha,dP)\,\mathcal U_{(0,1)}(dt)\,
 q_{t\mid1}^P(dx\mid\operatorname{Enc}_P(G,\alpha)).
\end{equation}
For fixed $P$, the encoded endpoint law $q_1^P$ induces
\begin{equation}
 q_t^P(x)=\sum_{x_1}q_1^P(x_1)q_{t\mid1}^P(x\mid x_1),\qquad
 \rho_i(a\mid x,t,P)=\Pr_\nu[X_{\star,i}=a\mid X_t=x,t,P].
 \label{eq:population-path}
\end{equation}
Full source support on the finite allowed alphabets makes $q_t^P(x)>0$ for every $x\in\Omega_P$
and $t<1$. At $t=0$ the path is $q_0^P$; at $t=1$ it is $q_1^P$.

\paragraph{Conditional and marginal rates.}
For $i\in\mathcal V(P)$ and $a,u,v\in\mathcal X_i(P)$, write
$s_{t,i}^P(u\mid a)=(1-t)q_{0,i}^P(u)+t\mathbf1\{u=a\}$ and
$\dot s_i^P(u\mid a)=\mathbf1\{u=a\}-q_{0,i}^P(u)$.
The conditional rate used in the ideal analysis is the $R^*$ construction of
\citet{qin2025defog}, restricted to the allowed alphabet:
\begin{equation}
 R^*_{t,i}(u,v\mid a,P)
 =\frac{[\dot s_i^P(v\mid a)-\dot s_i^P(u\mid a)]_+}
 {|\mathcal X_i(P)|\,s_{t,i}^P(u\mid a)},\qquad u\neq v,\quad 0\leq t<1.
 \label{eq:ideal-conditional-rate}
\end{equation}
Diagonal entries are the negative row sums. Denominators are positive for $t<1$; no value at $t=1$
is needed. Fixed coordinates have zero rates. For $x\neq y$ in $\Omega_P$, the joint conditional
and predicted rates are
\begin{align}
 R_t^*(x,y\mid x_1,P)
 &=\sum_{i\in\mathcal V(P)}\mathbf1\{y_{-i}=x_{-i}\}
 R^*_{t,i}(x_i,y_i\mid x_{1,i},P),\\
 R_t^\theta(x,y\mid P)
 &=\sum_{i\in\mathcal V(P)}\mathbf1\{y_{-i}=x_{-i}\}
 R_{t,i}^\theta(x_i,y_i\mid x,P).
 \label{eq:joint-local-rate}
\end{align}
Here $x_{-i}$ omits coordinate $i$, and Equation~\ref{eq:posterior-rate} defines the predicted
coordinate rate. Diagonal joint rates again make each row sum to zero. States differing at two or
more coordinates have zero instantaneous transition rate.

\begin{proposition}[Program-fixed invariance]
\label{prop:fixed-invariance}
Fix $P\in\mathcal P_{\rm adm}$ and an endpoint in $\Omega_P$. Corruption by
Equation~\ref{eq:conditional-path} stays in $\Omega_P$. The ideal CTMC started from $q_0^P$ stays
there for $t<1$ whenever its coordinate predictions are probability distributions on the allowed
alphabets. Numerical sampling in Equation~\ref{eq:numerical-sampling}, started from $q_0^P$
and using the allowed coordinate alphabets, satisfies
$\Pr[Z_k\in\Omega_P\text{ for all }0\leq k\leq L_{\rm samp}]=1$.
\end{proposition}
\begin{proof}
The corruption law has zero mass outside $\Omega_P$, and the ideal generator has no transition to
its complement. Each numerical coordinate draw stays in its allowed alphabet and $\Pi_P$ restores
the prescribed fixed values. Induction on $k$ proves the discrete statement. These arguments do not
require an accurate denoiser and do not imply molecular validity.
\end{proof}

\begin{proposition}[Posterior minimizer of the reference loss]
\label{prop:masked-bayes}
Under the reference training measure $\nu$, assume unrestricted measurable coordinate predictors
and positive weights $w_i(P)$ for each variable coordinate. Every population minimizer of
Equation~\ref{eq:primary-loss} satisfies
$d_{\theta,i}(a\mid x,t,P)=\rho_i(a\mid x,t,P)$ for all $a\in\mathcal X_i(P)$,
$\nu$-almost surely on contexts with $i\in\mathcal V(P)$.
\end{proposition}
\begin{proof}
Condition on $(X_t,t,P)$ and a variable coordinate. Since $w_i(P)$ is positive and does not depend
on its clean target after this conditioning, its conditional loss is
\begin{equation}
 w_i(P)\left[H(\rho_i)+D_{\rm KL}(\rho_i\Vert d_{\theta,i})\right].
\end{equation}
The entropy is independent of the predictor, and the divergence is nonnegative with equality only
at $d_{\theta,i}=\rho_i$. An unrestricted class can attain all these pointwise minima together.
A population minimizer must attain them almost surely, since a positive-measure excess would
increase the integrated loss. Fixed coordinates contribute no term.
\end{proof}

\begin{corollary}[Population recovery of the local rate]
\label{cor:local-rate-recovery}
For the reference population minimizer in Proposition~\ref{prop:masked-bayes},
$R_t^{\theta^\star}(x,y\mid P)$ equals the conditional expectation of
$R_t^*(x,y\mid X_\star,P)$ given $(X_t=x,t,P)$, almost surely under the path measure.
\end{corollary}
\begin{proof}
Each local conditional rate depends on the endpoint only through its coordinate $X_{\star,i}$.
Substituting the posterior marginal into Equation~\ref{eq:posterior-rate} therefore computes the
required conditional expectation. Summing over coordinates gives the joint rate.
\end{proof}

\begin{proof}[Proof of Theorem~\ref{thm:ideal-terminal-law}]
We verify the forward equation first for one coordinate, then for the conditional product path,
and finally after averaging over endpoints. Fix $P$, $x_1$ and $i$; suppress these indices in
$s_t$ and $\dot s$. For each $v$ and $t<1$, the net flux of Equation~\ref{eq:ideal-conditional-rate}
is
\begin{align}
 \sum_u s_t(u)R_t^*(u,v)
 &=\frac{1}{|\mathcal X_i(P)|}\sum_u
 \bigl([\dot s(v)-\dot s(u)]_+-[\dot s(u)-\dot s(v)]_+\bigr)\\
 &=\dot s(v)-\frac{1}{|\mathcal X_i(P)|}\sum_u\dot s(u)
 =\dot s(v).
\end{align}
The last equality uses $\sum_u\dot s(u)=0$. The product rule then gives the conditional joint
forward equation because the generator changes one coordinate at a time. Averaging over $q_1^P$
and using Bayes' rule yields
\begin{align}
 \partial_t q_t^P(y)
 &=\sum_{x,x_1}q_1^P(x_1)q_{t\mid1}^P(x\mid x_1)R_t^*(x,y\mid x_1,P)\\
 &=\sum_x q_t^P(x)\,
 \mathbb E[R_t^*(x,y\mid X_\star,P)\mid X_t=x,t,P].
\end{align}
Corollary~\ref{cor:local-rate-recovery} identifies the last conditional expectation with the
learned ideal rate for almost every $t$. On each interval $[0,1-\eta]$, $0<\eta<1$, source
positivity bounds all denominators away from zero. The rates are bounded measurable functions on a
finite state space, so the forward equation has a unique solution for initial law $q_0^P$.
Since $q_t^P$ is positive throughout $\Omega_P$, almost-sure equality under the path measure implies
rate equality at every state for Lebesgue-almost every $t$. Changes on the exceptional times do not
change the integral forward equation. The exactly simulated process therefore has law $q_t^P$.

The finite product path converges to the endpoint point mass as $t\uparrow1$. Summing over its
finite endpoint support gives $q_t^P\to q_1^P$ in total variation. This establishes a limit of
marginal laws; a pathwise value at the singular endpoint is not required. Exact decode round trip
maps each endpoint to its source constitution, proving
$\operatorname{Dec}_{\#}q_1^P=\mu_G^P$. The assumed endpoint validity and verified L1 give the
probability-one conclusion. If there are no variable coordinates, the compatible state space is a
singleton and the same conclusions hold for the constant process.
\end{proof}

\begin{corollary}[Decoder contraction and failure probability]
\label{cor:decoder-contraction}
Let $q_1^P$ satisfy the endpoint assumptions of Theorem~\ref{thm:ideal-terminal-law}. Let
$\widetilde q^P$ be any law on
$\overline{\mathfrak X}_{N,K_{\rm cl}}\cup\{\bot_{\rm term}\}$, with $q_1^P$ extended by zero
mass at $\bot_{\rm term}$ and $\operatorname{Dec}(\bot_{\rm term})=\bot_{\rm dec}$.
Define $\mathcal B_P$ as states whose decode fails, is molecularly invalid or fails verified L1.
Then
\begin{equation}
 \|\operatorname{Dec}_{\#}\widetilde q^P-\mu_G^P\|_{\rm TV}
 \leq\|\widetilde q^P-q_1^P\|_{\rm TV},\qquad
 \widetilde q^P(\mathcal B_P)\leq\|\widetilde q^P-q_1^P\|_{\rm TV},
\end{equation}
where $\|p-q\|_{\rm TV}=\frac12\sum_x|p(x)-q(x)|$.
\end{corollary}
\begin{proof}
Every decoder-output event has a preimage in the raw space, so its probability difference is bounded
by the raw-state total variation. Taking the supremum over output events proves the first bound.
The endpoint assumptions give $q_1^P(\mathcal B_P)=0$; applying the same event bound to
$\mathcal B_P$ proves the second.
\end{proof}
The corollary can be applied to the law \emph{after} numerical sampling and terminal readout.
It does not estimate that law's distance from $q_1^P$. In particular, error before the
learned terminal readout cannot be substituted without controlling the readout's effect.

\subsection{Accepted counts under a fixed sampling budget}

Fix a trained model, requested chemistry $c$ and its program-sampling policy. An attempt outcome $Y$
contains its sampled program $P$ and decoded graph or typed failure. Let $Q$ be the common law of
these outcomes and define
\begin{equation}
 A(Y)=\mathbf1\{G\text{ is valid and }\widehat{\mathrm{L1}}(G,c)=1\},\qquad
 z_Q=\mathbb E_Q[A(Y)],
\end{equation}
with $A(Y)=0$ for failures. This definition permits a fresh morphology-specific $P$ on every
attempt. A comparison with a single fixed $P$ is the corresponding special case.

\begin{proposition}[Exact law of matched-budget post-hoc filtering]
\label{prop:posthoc-budget}
For an integer $B\geq0$ and independent attempts with common law $Q$, the accepted-attempt count
$N_B=\sum_{b=1}^BA(Y_b)$ is binomial with parameters $(B,z_Q)$. For all $z_Q\in[0,1]$,
\begin{equation}
 \mathbb E[N_B]=Bz_Q,\qquad
 \operatorname{Var}(N_B)=Bz_Q(1-z_Q),\qquad
 \Pr[N_B=0]=(1-z_Q)^B.
\end{equation}
For $B=0$ the last probability is one, including when $z_Q=1$. If $z_Q>0$, the conditional law of
an accepted attempt is $Q_{\rm post}(y)=Q(y)A(y)/z_Q$, and the number $T_m$ of attempts needed for
$m\geq1$ acceptances satisfies $\mathbb E[T_m]=m/z_Q$.
For two proposal laws under the same attempt budget,
$\mathbb E[N_B^{(1)}-N_B^{(0)}]=B(z_{Q_1}-z_{Q_0})$.
\end{proposition}
\begin{proof}
Each acceptance indicator is Bernoulli with parameter $z_Q$, and independence makes their sum
binomial. Its expectation and variance give the first two identities; all $B$ attempts fail with
probability $(1-z_Q)^B$. For $z_Q>0$, conditioning on acceptance gives the normalized law by
Bayes' rule. The waiting times between acceptances are independent geometric variables on
$\{1,2,\ldots\}$ with mean $1/z_Q$; summing $m$ of them proves the waiting-time formula.
Subtracting the expected counts proves the comparison identity, without requiring independence
between the two arms.
\end{proof}
If $z_Q=0$, accepted outcomes have no conditional law and $T_m=\infty$ almost surely for $m\geq1$.
The proposition counts accepted attempts, including duplicates. It does not establish that program
conditioning increases acceptance; that is an empirical comparison. If requested programs are fixed
in unequal strata, independent acceptance indicators generally have different probabilities and
their sum is Poisson-binomial. Retraining variation is separate from both sampling laws.

\subsection{Route termination and evidence completeness}

\begin{lemma}[Termination of bounded route assessment]
\label{lem:planner-termination}
Fix a route contract $\mathfrak B$ with finite maximum depth $D_{\mathfrak B}$. Assume finitely many
proposed expansions per molecule and reactants per expansion, a terminating evidence and
forward-verification call for each proposal, and a search that assesses each proposed expansion
finitely many times. For an admitted root trace, recursive assessment terminates with a dossier
or typed failure in $\mathfrak D\cup\mathfrak F_{\mathfrak B}$.
\end{lemma}
\begin{proof}
Induct on remaining depth. At depth zero, the terminating availability query either admits a
terminal or returns a failure without further expansion. At positive remaining depth, only finitely
many proposals are assessed, each using terminating local calls. Every proposed child has one
fewer remaining level, so its assessment terminates by induction. A finite number of finite child
assessments and proposals therefore terminates at the parent. The same argument applies if a
finite outer root loop is used. Rejecting ancestor identities prevents circular dependencies from
being admitted; the finite depth and no-infinite-retry assumptions establish termination.
\end{proof}

\begin{invariant}[Synthesis-dossier admission]
\label{inv:dossier-soundness}
Under the frozen contract $\mathfrak B$, a successful assessor return $T\in\mathfrak D$ is admitted
only after checking $\mathrm{Complete}_{\mathfrak B}(G,T,P)=1$. Every other outcome is recorded as
a typed failure. This is the admission rule in Algorithm~\ref{alg:forge-inference}, not a theorem
that bounded search finds every supported chemical route.
\end{invariant}
The predicate checks molecular validity, a unique verified root trace, a well-formed finite tree,
evidence admission and forward replay for every preparation step, and availability for every leaf.
Its factors are defined in Appendix~\ref{app:route-assessment}. A failure establishes only that the
assessor did not return a complete dossier under the recorded contract. It does not exclude a route
under a larger budget, other admitted chemistry or another supplier snapshot.

\clearpage

\section{Historical reaction schemes and generated examples}
\label{app:chemistry-examples}
\label{app:preliminaries}

The assembly schemes depict the chemistry introduced in Section~\ref{sec:reaction-background}.
The generated examples pair molecular structures with precursor origins and exact assembly traces.
Their evidentiary scope is structural and computational.

\subsection{Assembly schemes}
\label{app:chemistry}
The precursor roles and assembly patterns are summarized in Table~\ref{tab:reaction-family-background}.

\begin{table}[htbp]
\caption{\textbf{Assembly chemistry underlying the three reaction programs.}
The named roles specify reactive functionality without selecting a precursor identity. Repeated
programs encode the successive additions admitted by the corresponding registered transform.}
\label{tab:reaction-family-background}
\centering
\small
\begin{tabular}{@{}>{\raggedright\arraybackslash}p{0.18\textwidth}>{\raggedright\arraybackslash}p{0.30\textwidth}>{\raggedright\arraybackslash}p{0.46\textwidth}@{}}
\toprule
Reaction family & Role-assigned precursors & Assembly pattern \\
\midrule
Ugi 3-CR & Amine, aldehyde and isocyanide &
Forms an $\alpha$-amino amide combining three precursor contributions \citep{xu2024agile}. \\
\addlinespace
Repeated aza-Michael addition & Amine head and acrylate-bearing hydrophobic components &
Amine addition to acrylate groups forms C--N bonds; successive additions attach hydrophobic
components at available amine sites \citep{li2023combinatorial}. \\
\addlinespace
Repeated reductive amination & Amine head and aldehyde-bearing hydrophobic components &
Condensation followed by reduction forms C--N linkages, attaching hydrophobic components through
successive transformations \citep{xue2024barcoding}. \\
\bottomrule
\end{tabular}
\end{table}

The three registered assembly patterns use skeletal notation and explicit hydrogens
to mark substitution constraints at the reaction core (Figure~\ref{fig:reaction-schemes}).
Precursor-origin boundaries can differ from the functional headgroup, linker and tail regions
introduced in Section~\ref{sec:lipid-background}.

% [tbp] rather than [H]: the schemes are too tall to sit inline, and forcing them here broke the
% page early and left LaTeX stretching the section glue above to fill it.
\begin{figure}[tbp]
\centering
\input{figures/forge_reaction_schemes/schemes_body.tex}
\caption{\textbf{The three assembly chemistries.} Monochrome skeletal drawings show the
reactants and products of each registered assembly. Hydrogens are drawn only on the atoms whose substitution decides
whether a product can be decomposed back into precursors: the head N--H, and the
$\alpha$-carbon C--H that the Ugi reaction takes from the aldehyde. $R^i$ stand for the
source-linked components of the frozen library, which are larger than the fragments drawn here.
The Ugi variant has no carboxylic-acid reactant component.}
\label{fig:reaction-schemes}
\end{figure}

\subsubsection{Reaction transforms and exact assembly}

Each assembly chemistry defines a forward transformation that maps role-assigned precursor
molecules to the product graphs permitted by that reaction.
The corresponding reverse transformation decomposes a complete lipid into the precursor
assignments that could have produced it.
To verify a decomposition, we computationally replay the reaction on the recovered precursors.
The product passes the level-one (L1) assembly check only if the reconstruction matches the original
lipid exactly.

\clearpage

\subsection{Generated products and exact precursors}
\label{app:generated-examples}
Figure~\ref{fig:forge-generated-samples} pairs two generated Ugi products with their unique
verified exact-L1 precursors. The examples illustrate recovered component novelty within the
frozen population evaluated in Section~\ref{sec:results}.

\begin{figure}[htbp]
\caption{\textbf{Exact Ugi products with recovered precursors absent from training.}
Each row contains a deterministic seed-0 product with a unique verified exact-L1 decomposition and a
verifier-recovered precursor absent from training. Columns show the generated lipid and its recovered
precursors. These are
structural examples, not evidence of activity or experimental synthesis success.}
\label{fig:forge-generated-samples}
\centering
% Vector drawings use the unchanged frozen sample identities; see the chemical_drawings receipt.
\newlength{\samplerowheight}\setlength{\samplerowheight}{4.2cm}
\newlength{\sampleproductwidth}\setlength{\sampleproductwidth}{0.465\textwidth}
\newlength{\sampleprecursorwidth}\setlength{\sampleprecursorwidth}{0.49\textwidth}
\newcommand{\iclrsamplerow}[1]{%
\begin{minipage}[c][\samplerowheight][c]{\sampleproductwidth}\centering
\includegraphics[width=\linewidth]{figures/chemical_drawings/sample_#1_product.pdf}
\end{minipage}
&
\begin{minipage}[c][\samplerowheight][c]{\sampleprecursorwidth}\centering
\includegraphics[width=\linewidth]{figures/chemical_drawings/sample_#1_precursors.pdf}
\end{minipage}}
\setlength{\tabcolsep}{0pt}
\resizebox{\textwidth}{!}{%
\begin{tabular}{@{}c@{\hspace{2.6mm}}!{\color{atlasrule}\vrule width 0.5pt}@{\hspace{2.6mm}}c@{}}
\toprule
\atlashead{\sampleproductwidth}{Generated lipid}
& \atlashead{\sampleprecursorwidth}{Exact L1 building blocks}\\
\midrule
\addlinespace[1.2mm]
\iclrsamplerow{01}\\
\addlinespace[2.2mm]
\iclrsamplerow{02}\\
\addlinespace[0.6mm]
\bottomrule
\end{tabular}
}
\end{figure}
\FloatBarrier

\subsection{Generated-structure atlas}
\label{app:atlas}

Table~\ref{tab:forge-generated-atlas} reports 12 deterministic display examples
from the same frozen seed-0 Ugi population used in Figure~\ref{fig:forge-generated-samples}. We selected the product with the lowest constitutional SHA-256 rank first, then applied ECFP4
MaxMin selection among the 23 eligible, constitutionally unique products. Selection favored
structural diversity without using activity, route closure, procurement evidence or human preference. Each row includes the canonical
constitutional SMILES so that the displayed identity is explicit; the exact strings are also stored
in the hash-pinned atlas receipt.

\begingroup
\arrayrulecolor{atlasrule}
\setlength{\arrayrulewidth}{0.5pt}
\setlength{\LTcapwidth}{\textwidth}
\setlength{\LTleft}{0pt}
\setlength{\LTright}{0pt}
\setlength{\LTpre}{10pt}
\setlength{\LTpost}{0pt}
\begin{longtable}{@{}c@{\hspace{1.8mm}}c@{\hspace{2.8mm}}!{\color{atlasrule}\vrule width 0.5pt}@{\hspace{2.8mm}}c@{}}
\caption{\textbf{FORGE generated-structure atlas.}
The 12 rows form one multipage table, indexed by display order. Each generated constitutional
graph and its canonical SMILES are paired with the amine, aldehyde and isocyanide building blocks
recovered through exact atom-mapped forward replay. Exact L1 replay and descriptor-manifold
membership are structural diagnostics, not evidence of activity, route closure or synthesis success.
These are deterministic display examples, not experimental candidates.}
\label{tab:forge-generated-atlas}\\[2mm]
\toprule
& \atlashead{\atlasgraphwidth}{Generated lipid and canonical SMILES}
& \atlashead{\atlasprecursorwidth}{Exact L1 building blocks}\\
\midrule
\addlinespace[3mm]
\endfirsthead
\multicolumn{3}{@{}l}{\small\itshape Table~\ref{tab:forge-generated-atlas}, continued}\\[2mm]
\toprule
& \atlashead{\atlasgraphwidth}{Generated lipid and canonical SMILES}
& \atlashead{\atlasprecursorwidth}{Exact L1 building blocks}\\
\midrule
\addlinespace[3mm]
\endhead
\addlinespace[1mm]
\bottomrule
\endfoot
\addlinespace[1mm]
\bottomrule
\endlastfoot
\input{figures/forge_generated_sample_atlas_core_saturation_v1/atlas_rows.tex}%
\end{longtable}
\endgroup

\clearpage
\input{sections/historical_three_family_results.tex}

\end{document}
